\ifdefined\pdfsuppressptexinfo
  \pdfsuppressptexinfo=15
\fi
\documentclass[11pt]{article}
\pdfinfoomitdate=1
\pdftrailerid{}

\usepackage[T1]{fontenc}
\usepackage[utf8]{inputenc}
\usepackage{lmodern}
\usepackage[margin=1in]{geometry}
\usepackage{amsmath,amssymb,amsthm,mathtools}
\usepackage{booktabs}
\usepackage{array}
\usepackage{longtable}
\usepackage{enumitem}
\usepackage{microtype}
\usepackage{xcolor}
\usepackage{hyperref}

\hypersetup{
  colorlinks=true,
  linkcolor=blue!55!black,
  citecolor=blue!55!black,
  urlcolor=blue!55!black,
  pdftitle={Spectral chirality of lens spaces: level-wise helicity balance, eta congruences, and silver zeros for transverse-traceless tensor curl},
  pdfauthor={Luqman Mahmood}
}

\newtheorem{theorem}{Theorem}[section]
\newtheorem{proposition}[theorem]{Proposition}
\newtheorem{lemma}[theorem]{Lemma}
\newtheorem{corollary}[theorem]{Corollary}
\newtheorem{conjecture}[theorem]{Conjecture}
\theoremstyle{definition}
\newtheorem{definition}[theorem]{Definition}
\newtheorem{remark}[theorem]{Remark}
\newtheorem{example}[theorem]{Example}
\newtheorem{observation}[theorem]{Observation}

\DeclareMathOperator{\Sym}{Sym}
\DeclareMathOperator{\Tr}{Tr}
\DeclareMathOperator{\sgn}{sgn}
\DeclareMathOperator{\content}{content}
\newcommand{\TT}{\mathrm{TT}}
\newcommand{\Z}{\mathbb Z}
\newcommand{\Q}{\mathbb Q}
\newcommand{\R}{\mathbb R}
\newcommand{\C}{\mathbb C}
\newcommand{\ind}{\mathbf 1}
\newcommand{\e}{\mathrm e}
\newcommand{\cK}{\mathcal K}
\newcommand{\sawtooth}[1]{\left(\!\left(#1\right)\!\right)}
\newcolumntype{L}[1]{>{\raggedright\arraybackslash}p{#1}}

\title{Spectral chirality of lens spaces:\\
level-wise helicity balance, eta congruences, and silver zeros\\
for transverse-traceless tensor curl}
\author{Luqman Mahmood}
\date{1 October 2026\\[3pt]\small Preprint\\
DOI: \href{https://doi.org/10.5281/zenodo.23049613}{10.5281/zenodo.23049613}}

\begin{document}

\maketitle

\begin{abstract}
Let $A_J$ be the rank-$J$ projected transverse-traceless tensor-curl
operator on a round lens space $L(p,q)$, in the normalization of
\cite{MahmoodP1Collection2026}, with helicity branches
$V_{n+2J}\otimes V_n$ and $V_n\otimes V_{n+2J}$ and eta invariant
$\eta_J(L(p,q))$.  We study the finer question of whether the two
branches have equal multiplicity at every level, and the arithmetic of
$\eta_J$ that the answer leaves behind.  We write the level-wise helicity
imbalance as an explicit sum over group elements with a closed-form
trigonometric kernel, prove that balance at every level is equivalent to a
finite identity in the group ring $\Z[\Z/p]$ together with the absence of
noncentral Clifford-translation elements, and classify all balanced
triples $(p,q,J)$ with $p\geq3$: they are the mirror-symmetric spaces
$q^2\equiv-1\pmod p$, the ``silver'' family $p\mid4J^2-2$ with
$q\equiv\pm(2J\pm1)$, and the blind ranks $p\mid2J$ under the
Clifford condition.  At rank one this recovers, by composition of
classical results, the statement that the coexact curl spectrum of
$L(p,q)$ is symmetric level by level exactly when $L(p,q)$ is
amphicheiral.  We then prove a divisor decomposition of $p\,\eta_J$ into
primitive parts, the congruence $\eta_J-J(2J^2-1)\eta_1\in2\Z$ at every
modulus, the sharper prime-modulus congruence
$p\,\eta_J\equiv-J(2J^2-1)(q+q^{-1})/3\pmod p$ with $p\,\eta_J\in2\Z$
for primes $p\geq5$, and the resulting exact denominators.  These settle
the zero set of $\eta_J$ at prime moduli $p\geq5$ away from the silver
ranks $2J^2\equiv1\pmod p$: for $p\nmid J$ the zeros are exactly the
mirror spaces.  At silver ranks $\eta_J$ is an even integer, and we prove
that $\eta_J(L(p,J))=0$ for every divisor $p>1$ of $2J^2-1$, extending a
family of \cite{MahmoodP1Collection2026} from $p=2J^2-1$ to all its
divisors; the mechanism is the conjugacy invariance of Rademacher's
symbol $\Psi$ on a trace $-6$ matrix commuting with an integral matrix of
determinant $-1$.  Finally we record what a
vanishing Rademacher symbol does and does not detect: inert and
reciprocal hyperbolic classes force $\Psi=0$, neither is necessary, and
the matrices behind the silver zero lie in the inert class $R^2L^2$.
The zero set at the level $p=2J^2-1$ itself (silver $S$-units, a fifth
family, a Pell ladder and a period law) is the subject of a companion
paper \cite{SilverRankZeros2026}, whose main statements we summarize.
The complete zero set at silver ranks remains open.
\end{abstract}

\section{Introduction}

The paper \cite{MahmoodP1Collection2026} evaluates the eta invariant
$\eta_J(L(p,q))$ of the rank-$J$ projected transverse-traceless (TT)
tensor-curl operator on every round lens space, uniformly in the rank and
in prime and composite levels, and exhibits four infinite families of
zeros at $p=2J^2-1$.  It leaves open the classification of all zero
orbits ``beyond the four constructions above and the orientation-forced
condition $q^2\equiv-1\pmod p$''.  The eta invariant is a regularized sum
of the level-wise imbalance $D_J(n)=N^+_J(n)-N^-_J(n)$ between the two
helicity branches; a vanishing eta invariant may or may not come from a
vanishing imbalance at every level.  The present paper is a companion to
\cite{MahmoodP1Collection2026}.  It uses the same operator, the same
normalization and the same lens-space conventions, and it addresses three
questions in turn.

\begin{enumerate}[label=(\roman*)]
\item \emph{Level-wise balance.}  For which $(p,q,J)$ is $D_J(n)=0$ for
every level $n\geq0$?  Sections~\ref{sec:balance} and
\ref{sec:classification} answer this completely for $p\geq3$
(Theorems~\ref{thm:kernel}, \ref{thm:criterion} and
\ref{thm:classification}).  At rank one the answer is classical by
composition of three published results
\cite{Bar2019,Katase1987,RademacherGrosswald1972,PrzytyckiYasuhara2003,
Boldt2017}; for $J\geq2$ we found neither the criterion nor the families
in the literature we searched.
\item \emph{Congruences.}  Sections~\ref{sec:decomposition} and
\ref{sec:congruences} prove a decomposition of $p\,\eta_J$ over the
divisors of $p$, the congruence
$\eta_J-J(2J^2-1)\eta_1\in2\Z$ at every modulus (Theorem~\ref{thm:U}),
and at primes $p\geq5$ the finer statement
$p\,\eta_J\in2\Z$, $p\,\eta_J\equiv-J(2J^2-1)(q+q^{-1})/3\pmod p$
(Theorem~\ref{thm:E}).  These determine the exact denominator of
$\eta_J$ and settle the zero problem of \cite{MahmoodP1Collection2026}
at prime moduli $p\geq5$ away from the silver ranks $2J^2\equiv1\pmod p$.
\item \emph{Silver zeros.}  At a silver rank $\eta_J$ is an even integer
and the zero set is richer.  Section~\ref{sec:silver} proves
$\eta_J(L(p,J))=0$ for every divisor $p>1$ of $2J^2-1$
(Theorem~\ref{thm:C3}), the first family of
\cite{MahmoodP1Collection2026} extended from $p=2J^2-1$ to all its
divisors.  Section~\ref{sec:rademacher} explains the mechanism in the
language of Rademacher's symbol on hyperbolic conjugacy classes and
records, with citations and without novelty claims, what a vanishing
symbol detects and what it forgets.  Section~\ref{sec:silver-zero-set}
summarizes what the companion paper \cite{SilverRankZeros2026} proves
about the zero set at $p=2J^2-1$: the silver $S$-unit coordinate, a
fifth zero family, the Pell ladder, the period law for every orbit of
norm $\pm1$ or $\pm2$, and the unbounded ambiguity of a vanishing
$\eta_J$.
\end{enumerate}

Relation to the companion paper \cite{MahmoodP3Collection2026}: that
paper transfers twisted spin-Dirac eta vectors to all-rank TT eta vectors
on every free round spherical space form, and organizes the one-sided
quotients along affine-ADE diagrams.  It contains no level-wise
statements, no congruences and no zero classification, and it is not used
here; conversely nothing below concerns non-cyclic quotients.  The
boundary realization of the operator in \cite{MahmoodP7Collection2026}
is likewise not used.

Helicity-resolved tensor mode counts on lens spaces are also of
independent interest in cosmic topology, where the effect of topology on
spin-two eigenmodes was recently described as unstudied
\cite{COMPACTIIIa2024}; no physical claim is made in this paper.

Exact computations are recorded in Appendix~\ref{app:verification}.
They test the identities and conventions; they do not replace the
proofs of universal statements.

\section{Conventions}
\label{sec:conventions}

We keep the conventions of \cite[Sec.~2]{MahmoodP1Collection2026}
verbatim.  Write
$S^3=\{(z_1,z_2)\in\C^2:|z_1|^2+|z_2|^2=1\}$.  For $p\geq2$ and
$q\in(\Z/p\Z)^\times$ the generator of $\Z_p$ acts by
\[
 (z_1,z_2)\longmapsto\bigl(\e^{2\pi i/p}z_1,\e^{2\pi iq/p}z_2\bigr),
\]
and $L(p,q)$ is the quotient with its complex quotient orientation.  We
write $q^*$ for the inverse of $q$ modulo $p$, and $\tilde q$ for the odd
lift of $q$ ($\tilde q=q$ if $q$ is odd, $\tilde q=q+p$ if $q$ is even
and $p$ odd; if $p$ is even every unit is odd).  Put
\begin{equation}
 A=\frac{1+\tilde q}{2},\qquad B=\frac{1-\tilde q}{2},\qquad A+B=1.
 \label{eq:AB}
\end{equation}
Let $V_a$ be the irreducible $SU(2)$ representation of highest weight $a$
and dimension $a+1$, with weights $W_a=\{-a,-a+2,\ldots,a\}$.  The lifted
generator acts on the weight vector $(w_L,w_R)$ of $V_a\otimes V_b$ by
$\exp\bigl(2\pi i(Aw_L+Bw_R)/p\bigr)$; by
\cite[Lem.~5.2]{MahmoodP1Collection2026} (central descent) the invariant
counts below do not depend on the lift.  For $a\equiv b\pmod2$ this is the
cyclic invariant count $I_{a+1,b+1}(p,q)$ of
\cite[Sec.~2]{MahmoodP1Collection2026},
\begin{equation}
 N(a,b)=\#\{(w_L,w_R)\in W_a\times W_b:\ Aw_L+Bw_R\equiv0\pmod p\}
 =I_{a+1,b+1}(p,q).
 \label{eq:Nab}
\end{equation}

The rank-$J$ projected TT tensor-curl operator $A_J$ of
\cite[Thm.~3.1]{MahmoodP1Collection2026} has, on the unit sphere, the
signed spectrum
\[
 V_{n+2J}\otimes V_n:\ +(n+J+1),\qquad
 V_n\otimes V_{n+2J}:\ -(n+J+1),\qquad n\geq0,
\]
and on $L(p,q)$ the eigenspaces are the invariant vectors.  For a
\emph{level} $n\geq0$ put
\begin{equation}
 N^+_J(n)=N(n+2J,n),\qquad N^-_J(n)=N(n,n+2J),\qquad
 D_J(n)=D_J(n;p,q)=N^+_J(n)-N^-_J(n),
 \label{eq:branch-counts}
\end{equation}
and write $\nu=n+J+1$ for the eigenvalue magnitude.  The eta invariant is
\begin{equation}
 \eta_J(L(p,q))=\eta_{A_J}(0),\qquad
 \eta_{A_J}(s)=\sum_{n\geq0}\frac{D_J(n)}{(n+J+1)^s},
 \label{eq:eta-def}
\end{equation}
continued to $s=0$ as in \cite[Prop.~4.1]{MahmoodP1Collection2026}; the
continuation is regular at $s=0$ and its value equals the Abel value of
the corresponding root-of-unity exponential-polynomial series
\cite[Lem.~5.4]{MahmoodP1Collection2026}.  We use the lens symmetries
\begin{equation}
 \eta_J(L(p,q^*))=\eta_J(L(p,q)),\qquad
 \eta_J(L(p,-q))=-\eta_J(L(p,q)),
 \label{eq:lens-symmetries}
\end{equation}
the rank-one value $\eta_1(L(p,q))=-4s(q,p)$
\cite[Cor.~5.8]{MahmoodP1Collection2026}, where
\[
 \sawtooth{x}=
 \begin{cases}x-\lfloor x\rfloor-\tfrac12,&x\notin\Z,\\0,&x\in\Z,\end{cases}
 \qquad
 s(q,p)=\sum_{a=0}^{p-1}\sawtooth{\frac ap}\sawtooth{\frac{qa}p},
\]
and the all-rank master formula of \cite[Thm.~5.1]{MahmoodP1Collection2026},
\begin{align}
 \eta_J(L(p,q))={}&-4J\,s(q,p)
 +\frac2p\sum_{a=1}^{J-1}(J-a)\bigl[\rho_p(aq)+\rho_p(aq^*)\bigr]\nonumber\\
 &-2\!\!\sum_{\substack{a,b\geq1\\a+b\leq J-1}}\!\!(J-a-b)
 \bigl[\delta_p(aq-b)-\delta_p(aq+b)\bigr]
 -L_J(p,q),
 \label{eq:master}
\end{align}
with $\rho_p(k)=0$ if $p\mid k$ and $p-2\langle k\rangle_p$ otherwise,
$\delta_p(k)=\ind_{p\mid k}$, and the low-mode term
$L_J(p,q)=\sum_{m=1}^{J-1}[I_{J-m,J+m}(p,q)-I_{J+m,J-m}(p,q)]$.
Equation~\eqref{eq:master} holds for every $p\geq2$, every unit $q$ and
every $J\geq1$.  Its last two terms combine into a single signed sum
over lattice points, the box form of
\cite[Lem.~8.1]{SilverRankZeros2026}; that is a rewriting of these two
terms, not a new invariant, and nothing in this paper uses it.

\begin{definition}
\label{def:terms}
Let $p\geq1$, $q$ a unit modulo $p$ and $J\geq1$.
\begin{enumerate}[label=(\alph*)]
\item $L(p,q)$ is \emph{rank-$J$ balanced} (balanced at every level)
if $D_J(n;p,q)=0$ for all $n\geq0$.
\item $L(p,q)$ is \emph{mirror} (amphicheiral) if $q^2\equiv-1\pmod p$;
otherwise it is \emph{chiral}.  The condition is the classical criterion
for an orientation-reversing self-homeomorphism, indeed an
orientation-reversing isometry
\cite{PrzytyckiYasuhara2003,Boldt2017}: if $q^2\equiv-1$ then
$T(z_1,z_2)=(\bar z_2,z_1)$ normalizes the deck group and reverses
orientation.
\item The rank $J$ is a \emph{silver rank} for $p$ if $2J^2\equiv1\pmod p$
(equivalently $p\mid2J^2-1$, so $p$ is odd and coprime to $6$).
\item The rank $J$ is a \emph{blind rank} for $p$ if $p\mid2J$.
\item The \emph{Clifford condition} is
\begin{equation}
 \gcd(q-1,p)\leq2\quad\text{and}\quad\gcd(q+1,p)\leq2.
 \tag{C}
 \label{eq:clifford}
\end{equation}
For odd $p$ it says $\gcd(q^2-1,p)=1$; for even $p$ it says that both
gcds equal $2$.  It has nonempty scope only when $\gcd(p,12)\in\{1,2\}$
(Proposition~\ref{prop:existence}).
\end{enumerate}
\end{definition}

The mirror condition and the two lens symmetries
\eqref{eq:lens-symmetries} give the two ``forced'' facts that anchor
everything below: a mirror space has $D_J(n)=0$ at every level and every
rank, because the orientation-reversing isometry exchanges the two
branches; and $\eta_J(L(p,q))=0$ whenever $q^*\equiv-q$, by
\eqref{eq:lens-symmetries}.

\paragraph{Notation for group elements.}
For $0\leq k<p$ put $\theta_k=2\pi k/p$.  Let
\begin{equation}
 D^+=\{k:\ p\mid k(q-1)\},\qquad D^-=\{k:\ p\mid k(q+1)\},\qquad
 Z=\{k:\ 2k\equiv0\pmod p\},
 \label{eq:degenerate-sets}
\end{equation}
so that $Z=\{0\}$ or $\{0,p/2\}$, and put $D^{\pm*}=D^\pm\setminus Z$ and
$\mathcal N=\Z/p\setminus(D^+\cup D^-)$.  The element $g^k$ with
$k\in D^{+*}$ (resp.\ $D^{-*}$) acts on $S^3$ as a pure left (resp.\ right)
Clifford translation; we call these the \emph{Clifford-translation
elements}.  Condition~\eqref{eq:clifford} says exactly that
$D^{+*}=D^{-*}=\emptyset$, since $D^+$ is the subgroup of order
$\gcd(q-1,p)$ and $D^-$ the subgroup of order $\gcd(q+1,p)$.  Finally
$\chi_a(t)=\sin((a+1)t)/\sin t$ is the $SU(2)$ character
(with its removable limits $\chi_a(0)=a+1$, $\chi_a(\pi)=(-1)^a(a+1)$),
\begin{equation}
 c_J(\nu,t)=\frac{2\bigl[\nu\cos(\nu t)\sin(Jt)-J\sin(\nu t)\cos(Jt)\bigr]}
 {\sin t},
 \qquad
 \cK_J(t)=\frac{\sin(Jqt)}{\cos(qt)-\cos t}+\frac{\sin(Jq^*t)}{\cos(q^*t)-\cos t},
 \label{eq:cJ-KJ}
\end{equation}
and for $a\in\Z/p$ we write $[a]=\delta_a-\delta_{-a}\in\Z[\Z/p]$, so
that $[a]=0$ iff $2a\equiv0$.

\section{Level-wise helicity balance}
\label{sec:balance}

\subsection{The imbalance as a sum over group elements}

The per-element signed character of \cite[Sec.~5.1]{MahmoodP1Collection2026}
is, for a nonidentity element $g^k$ and $\nu=n+J+1$,
\begin{equation}
 E_k(n)=\chi_{n+2J}(x_k)\chi_n(y_k)-\chi_n(x_k)\chi_{n+2J}(y_k),
 \qquad
 x_k=\frac{2\pi kA}{p},\quad y_k=\frac{2\pi kB}{p},
 \label{eq:Ek}
\end{equation}
with $u_k:=x_k-y_k=2\pi k\tilde q/p\equiv q\theta_k$ and
$v_k:=x_k+y_k=\theta_k$ modulo $2\pi$.  Character averaging gives
$pD_J(n)=\sum_{k=0}^{p-1}E_k(n)$.  The following theorem reorganizes this
sum.  Its content is the closed-form kernel $\cK_J$ and the explicit
separation of the Clifford-translation elements; the per-element
character itself is from \cite{MahmoodP1Collection2026}, and at $J=1$ the
generating-function form of the same sum is in
\cite[Thm.~5.5, Cor.~5.8, Ex.~5.9]{Bar2019}.

\begin{theorem}[Kernel form of the level-wise imbalance]
\label{thm:kernel}
For $p\geq1$, every unit $q$, $J\geq1$ and $n\geq0$, with $\nu=n+J+1$,
\begin{equation}
 p\,D_J(n;p,q)=\sum_{k\in D^{+*}}c_J(\nu,\theta_k)
 -\sum_{k\in D^{-*}}c_J(\nu,\theta_k)
 +2\sum_{k\in\mathcal N}\sin(\nu\theta_k)\,\cK_J(\theta_k).
 \label{eq:kernel-form}
\end{equation}
The third sum is the generic part of $\sum_kE_k$ after the reindexing
$k\mapsto q^*k$ of the second half of each character difference; its
$k$-th summand is not in general the single-element contribution
$E_k(n)$.
\end{theorem}

\begin{proof}
\emph{Step 1.}  For $k\in Z$, $E_k=0$: at $k=0$ the two dimensions
agree, and at $k=p/2$ (even $p$) both $x_k$ and $y_k$ lie in $\pi\Z$, so
the character values are $(-1)^{\text{weight}}\times$ dimension and the
two products agree (\cite[Lem.~5.3]{MahmoodP1Collection2026}).

\emph{Step 2 (generic elements).}  Since
$\sin x\sin y=\tfrac12(\cos u-\cos v)$ with $u=x-y$, $v=x+y$, the
denominator of $E_k$ vanishes iff $u_k\equiv\pm v_k\pmod{2\pi}$ iff
$k\in D^+\cup D^-$.  For $k\in\mathcal N$, expanding the numerator by
product-to-sum with $a=\nu+J$, $b=\nu-J$ gives
\[
 \sin(ax)\sin(by)-\sin(bx)\sin(ay)
 =\sin(\nu v)\sin(Ju)-\sin(\nu u)\sin(Jv),
\]
hence
\begin{equation}
 E_k(n)=\frac{2\bigl[\sin(\nu v_k)\sin(Ju_k)-\sin(\nu u_k)\sin(Jv_k)\bigr]}
 {\cos u_k-\cos v_k}
 =\frac{2\sin(\nu\theta_k)\sin(Jq\theta_k)}{\cos(q\theta_k)-\cos\theta_k}
 -\frac{2\sin(\nu q\theta_k)\sin(J\theta_k)}{\cos(q\theta_k)-\cos\theta_k}.
 \label{eq:Ek-generic}
\end{equation}

\emph{Step 3 (reindexing).}  Multiplication by the unit $q$ preserves
$D^+$, $D^-$ and $\mathcal N$.  Substituting $k=q^*k'$ in the second term
of \eqref{eq:Ek-generic} and summing over $k\in\mathcal N$ gives
$\sum_{k'\in\mathcal N}2\sin(\nu\theta_{k'})\sin(Jq^*\theta_{k'})/
(\cos(q^*\theta_{k'})-\cos\theta_{k'})$.  Adding the first term yields
$2\sum_{\mathcal N}\sin(\nu\theta_k)\cK_J(\theta_k)$.

\emph{Step 4 (Clifford translations).}  For $k\in D^{+*}$ one has
$y_k\in\pi\Z$, say $y_k=\pi j$, and $\sin\theta_k\neq0$.  With
$\sigma=(-1)^j$, $\chi_a(y_k)=\sigma^a(a+1)$ and
$\chi_a(x_k)=\chi_a(\theta_k-\pi j)=\sigma^a\chi_a(\theta_k)$, so
\begin{align*}
 E_k(n)&=(n+1)\chi_{n+2J}(\theta_k)-(n+2J+1)\chi_n(\theta_k)\\
 &=\frac{(\nu-J)\sin((\nu+J)\theta_k)-(\nu+J)\sin((\nu-J)\theta_k)}{\sin\theta_k}
 =c_J(\nu,\theta_k).
\end{align*}
For $k\in D^{-*}$, $x_k\in\pi\Z$ and symmetrically $E_k(n)=-c_J(\nu,\theta_k)$.
Since $D^+\cap D^-\subseteq Z$, the three sums in \eqref{eq:kernel-form}
run over disjoint sets.
\end{proof}

\begin{corollary}[Affine behaviour on residue classes; finite window]
\label{cor:affine}
Write $p\,D_J(n)=\nu F(\nu)+G(\nu)$ with
\begin{equation}
 F(\nu)=\sum_{k}\varphi(k)\cos(\nu\theta_k),\qquad
 G(\nu)=\sum_{k}\psi(k)\sin(\nu\theta_k),
 \label{eq:FG}
\end{equation}
where $\varphi=\dfrac{2\sin(J\theta)}{\sin\theta}(\ind_{D^{+*}}-\ind_{D^{-*}})$
is even and
$\psi=-\dfrac{2J\cos(J\theta)}{\sin\theta}(\ind_{D^{+*}}-\ind_{D^{-*}})+2\cK_J\ind_{\mathcal N}$
is odd in $k$.  Both $F$ and $G$ are $p$-periodic in $\nu$.  Consequently
for each $0\leq r<p$ there are integers $u_r,v_r$ with
\begin{equation}
 D_J(r+pk;p,q)=u_rk+v_r\qquad(k\geq0),
 \label{eq:affine-mod-p}
\end{equation}
and $L(p,q)$ is rank-$J$ balanced iff $D_J(n)=0$ for $0\leq n<2p$.
\end{corollary}

\begin{proof}
Read off from \eqref{eq:kernel-form} and \eqref{eq:cJ-KJ}.  Two values of
$k$ determine an affine function.
\end{proof}

The affine statement on classes modulo $2p$, which already gives the
window $n<4p$, is \cite[Prop.~4.1]{MahmoodP1Collection2026}; the sharper
modulus $p$ can also be obtained directly by an elementary character
argument (the increment $c_{h+p}-c_h$ of the residue-class weight counts
depends only on $h$ modulo $2$).

\subsection{The criterion}

\begin{theorem}[Complete criterion for level-wise balance]
\label{thm:criterion}
Let $p\geq3$, $q$ a unit modulo $p$ and $J\geq1$.  Then $L(p,q)$ is
rank-$J$ balanced if and only if both
\begin{enumerate}[label=(\roman*)]
\item the Clifford condition \eqref{eq:clifford} holds, i.e.\
$\gcd(q-1,p)\leq2$ and $\gcd(q+1,p)\leq2$; and
\item in $\Z[\Z/p]$, with $X=Jq$ and $Y=Jq^*$,
\begin{equation}
 [X+q^*]+[X-q^*]+[Y+q]+[Y-q]=[X+1]+[X-1]+[Y+1]+[Y-1].
 \label{eq:odd-measure}
\end{equation}
\end{enumerate}
For $p\in\{1,2\}$ every $L(p,q)$ is balanced at every rank.
\end{theorem}

\begin{proof}
\emph{Necessity.}  Suppose $D_J(n)=0$ for all $n\geq0$.  Each residue
class of $\nu$ modulo $p$ contains infinitely many $\nu\geq J+1$, and
$\nu F(\nu)+G(\nu)$ is affine on the class, so $F=G=0$ on $\Z/p$.  Since
$\varphi$ is even and $\psi$ is odd, $F$ is the full Fourier transform of
$\varphi$ and $G$ that of $\psi$ (up to the factor $i$), so Fourier
inversion gives $\varphi=\psi=0$.
\begin{itemize}
\item $\varphi=0$ gives $\sin(J\theta_k)=0$ on $D^{+*}\cup D^{-*}$
(the two sets are disjoint, as $k$ in both forces $p\mid2k$), hence
$|\cos(J\theta_k)|=1$ there.
\item Then $\psi=\mp2J\cos(J\theta_k)/\sin\theta_k\neq0$ on $D^{\pm*}$,
contradicting $\psi=0$ unless $D^{+*}=D^{-*}=\emptyset$.  Since $D^\pm$ is
the subgroup of order $\gcd(q\mp1,p)$, this is (i).
\item Finally $\psi=0$ on $\mathcal N$ says $\cK_J(\theta_k)=0$ for
$k\in\mathcal N$, which is (ii) by the lemma below.
\end{itemize}
\emph{Sufficiency.}  Under (i), $D^{\pm*}=\emptyset$ and
$pD_J(n)=2\sum_{\mathcal N}\sin(\nu\theta_k)\cK_J(\theta_k)$, which
vanishes identically when $\cK_J=0$ on $\mathcal N$, i.e.\ under (ii).

\emph{Lemma.}  Under (i), $\cK_J=0$ on $\mathcal N$ iff \eqref{eq:odd-measure}.
Indeed, on $\mathcal N$ both denominators of $\cK_J$ are nonzero, so
$\cK_J(\theta)=0$ iff
\[
 S(\theta):=\sin(X\theta)\bigl(\cos(q^*\theta)-\cos\theta\bigr)
 +\sin(Y\theta)\bigl(\cos(q\theta)-\cos\theta\bigr)=0 .
\]
Product-to-sum gives
$2S(\theta)=\sum_{a\in P}\sin(a\theta)-\sum_{b\in M}\sin(b\theta)$ with
$P=(X+q^*,X-q^*,Y+q,Y-q)$ and $M=(X+1,X-1,Y+1,Y-1)$.  At $k\in Z$
(where $\theta\in\{0,\pi\}$) $S$ vanishes automatically, and under (i)
$\mathcal N\cup Z=\Z/p$; so $\cK_J=0$ on $\mathcal N$ iff $S(\theta_k)=0$
for all $k\in\Z/p$.  Now $S$ is the sine transform of the odd measure
$\mu=\sum_P[a]-\sum_M[b]$: $\sum_r\mu(r)\e^{ir\theta}=4iS(\theta)$.
By Fourier inversion $S\equiv0$ iff $\mu=0$, which is \eqref{eq:odd-measure}.

For $p=1$ the two branches have equal dimension; for $p=2$ both branch
weights have the parity of $n$ and the invariant counts agree.
\end{proof}

\subsection{Consequences and the three sufficient families}

\begin{corollary}
\label{cor:criterion-consequences}
Let $p\geq3$.
\begin{enumerate}[label=(\alph*)]
\item \textup{(Descent.)}  If $L(p,q)$ is rank-$J$ balanced, so is
$L(d,q\bmod d)$ for every $d\mid p$.
\item \textup{(Rank periodicity and reflection, level-wise.)}  The
rank-$J$ verdict depends on $J$ only modulo $p$, and $J$ and $-J$ give
the same verdict.
\item \textup{(Mirror.)}  If $q^2\equiv-1\pmod p$, then $L(p,q)$ is
rank-$J$ balanced for every $J$.
\item \textup{(Silver family.)}  If $p\mid4J^2-2$ and
$q\equiv\pm(2J-1)$ or $q\equiv\pm(2J+1)\pmod p$, then $L(p,q)$ is
rank-$J$ balanced and chiral.
\item \textup{(Blind ranks.)}  If $p\mid2J$ and \eqref{eq:clifford} holds,
then $L(p,q)$ is rank-$J$ balanced.
\item If $L(p,q)$ is rank-$J$ balanced, then $\eta_J(L(p,q))=0$.
\end{enumerate}
\end{corollary}

\begin{proof}
(a) Condition (i) passes to divisors, and $S=0$ on $\Z/p$ restricts to
$S=0$ on the subgroup $(p/d)\Z/p$, which is the function $S$ of
$L(d,q\bmod d)$.
(b) \eqref{eq:odd-measure} involves $J$ only through $X,Y$ modulo $p$,
and $J\mapsto-J$ sends $\mu\mapsto-\mu$.
(c) If $d\mid\gcd(q\mp1,p)$ then $q^2+1\equiv2\pmod d$, so $d\leq2$, and
(i) holds; moreover $q^*\equiv-q$, so $Y\equiv-X$ and the four $P$-terms
cancel in pairs, as do the $M$-terms.  (Alternatively, the
orientation-reversing isometry exchanges the branches.)
(d) The verdict is invariant under $q\mapsto-q$ (which negates $\mu$) and
$q\mapsto q^*$ (which exchanges $X$ and $Y$), so take $q=2J-1$; then
$q^*=2J+1$ because $(2J-1)(2J+1)=2(2J^2-1)+1$.  Put $e=2J^2-1$ (odd) and
$c=0$ if $p$ is odd, $c=p/2$ if $p$ is even.  If $p$ is even then $p=2d$
with $d\mid e$ odd and $e\equiv d\pmod{2d}$; in both cases
$X=J(2J-1)=e-J+1\equiv c+1-J$ and $Y=J(2J+1)\equiv c+1+J$.  Hence
$P=(c+J+2,\ c-3J,\ c+3J,\ c+2-J)$ and $M=(c+2-J,\ c-J,\ c+J+2,\ c+J)$.
Using $[c-a]=-[c+a]$ (as $2c\equiv0$), the terms $[c-3J]+[c+3J]$ and
$[c-J]+[c+J]$ vanish and the remaining terms match, proving
\eqref{eq:odd-measure}.  Condition (i): $\gcd(2J-2,2e)=\gcd(2J,2e)=2$
because $e=2J\cdot J-1$ and $e=2(J-1)(J+1)+1$.  Chirality:
$(2J-1)^2+1\equiv-4(J-1)\pmod p$, and $p\mid4(J-1)$ is impossible for
$p\geq3$ since the odd part of $p$ divides $e$ and $\gcd(e,J-1)=1$.
(e) $2X\equiv2Y\equiv0$ gives $[X+a]+[X-a]=0$ for every $a$, so both
sides of \eqref{eq:odd-measure} vanish.
(f) The eta function \eqref{eq:eta-def} of a balanced space vanishes
identically.
\end{proof}

\begin{remark}
\label{rem:family-b-versus-paper1}
Family (d) contains, for $p=2J^2-1$, the zeros $q=2J-1$ of
\cite[Thm.~7.1]{MahmoodP1Collection2026}, which form a uniformly
balanced family (balanced at every $J$); the other three
parametrizations of that theorem are not balanced in general, although
individual members can lie on the balanced or mirror loci (see
Example~\ref{ex:zero-not-balanced}).  It also contains every proper
divisor $p$ of $2J^2-1$ and the even levels $p=2d$, $d\mid2J^2-1$, which
lie outside that theorem (``every level $p_J$ in these families is
odd'').  Some proper-divisor members are reachable from
\cite{MahmoodP1Collection2026} by combining its zero families with its
rank symmetries (for instance $\eta_5(L(7,q))=-\eta_2(L(7,q))=0$ for
$q\in\{2,3,4,5\}$); others
are not, for instance $L(23,4)$ at $J=9$, where $23\mid2\cdot81-1=7\cdot23$,
no $J'$ has $2J'^2-1=23$, and the rank symmetries only relate
$J\equiv\pm9\pmod{23}$.  Every even level $p\equiv2\pmod4$ lies outside
both \cite[Thm.~7.1]{MahmoodP1Collection2026} (odd levels) and
\cite[Thm.~6.2]{MahmoodP1Collection2026} (which needs $\gcd(p,6)=1$).
\end{remark}

\begin{remark}[Mechanism of the blind ranks]
If $p\mid2J$ then $\sin(Jq\theta_k)=\sin(Jq^*\theta_k)=0$ for every $k$,
so the kernel $\cK_J$ vanishes at every class.  This is the same vanishing
of the $SU(2)$ character $\chi_{J-1}$ at every noncentral class that
underlies the blind adjacent-rank families of
\cite{MahmoodP7Collection2026}, and condition \eqref{eq:clifford} is the
primality-free non-centrality condition used there.  What is new at the
level-wise scale is the \emph{necessity} of \eqref{eq:clifford} for
balance, the criterion \eqref{eq:odd-measure}, and the silver family (d).
\end{remark}

\subsection{Boundedness, existence, and even levels}

\begin{proposition}[Bounded versus growing imbalance]
\label{prop:bounded}
Let $p\geq3$.  The sequence $D_J(n;p,q)$ is bounded in $n$ iff
$\gcd(q-1,p)\mid2J$ and $\gcd(q+1,p)\mid2J$.  In that case $D_J$ is
$p$-periodic in $n$; otherwise $|D_J|$ grows linearly along at least one
residue class modulo $p$.  A Clifford-translation element (a $k\in D^{+*}\cup D^{-*}$)
precludes balance at every level, but need not force growth: for
$L(3,1)$ and $J=3$ one has $D_3=(-4,4,0)$ repeated with period $3$, with
$(N^+,N^-)=(3,7),(4,0),(9,9)$ at $n=0,1,2$.
\end{proposition}

\begin{proof}
By \eqref{eq:FG}, $D_J$ is bounded iff $F\equiv0$ iff $\varphi=0$ iff
$\sin(J\theta_k)=0$ on every noncentral element of $D^+$ and of $D^-$.
On the subgroup $D^+$ of order $d=\gcd(q-1,p)$ this holds iff $d\mid2J$,
and it is automatic for $d\leq2$; likewise for $D^-$.  If $F=0$ then
$pD_J=G$ is $p$-periodic; if $F\neq0$ some class has nonzero slope.  The
last assertion is Theorem~\ref{thm:criterion}: with $F=0$ the residual
$G$ detects the Clifford elements.
\end{proof}

\begin{proposition}[Existence of balanced ranks]
\label{prop:existence}
For $p\geq3$, some pair $(q,J)$ with $L(p,q)$ rank-$J$ balanced exists
iff $\gcd(p,12)\in\{1,2\}$.  A chiral balanced example exists iff, in
addition, $p\notin\{5,10\}$.
\end{proposition}

\begin{proof}
If $3\mid p$, every unit $q$ is $\equiv\pm1\pmod3$, so
\eqref{eq:clifford} fails; if $4\mid p$, $q$ is odd and $4$ divides one of
$q\mp1$, so \eqref{eq:clifford} fails.  Conversely, for odd $p$ coprime
to $6$ take $q=2$, $J=p$; for $p=2d$ with $d$ odd and $3\nmid d$ take
$q=3$, $J=d$: both satisfy \eqref{eq:clifford} and $p\mid2J$, so
Corollary~\ref{cor:criterion-consequences}(e) applies.  For chirality,
$q=2$ is mirror only when $p=5$ and $q=3$ only when $p=10$; at $p=5$ the
admissible $q$ are $2,3$ and at $p=10$ they are $3,7$, all mirror.
\end{proof}

\begin{lemma}[Even levels are a parity projection of their odd parent]
\label{lem:parity-projection}
Let $p$ be odd and $\tilde q$ the odd lift of $q$.  Then for every $J$
and every level $n$,
\[
 N^\pm_{J}(n;2p,\tilde q)=
 \begin{cases}N^\pm_J(n;p,q),&n\text{ even},\\0,&n\text{ odd}.\end{cases}
\]
Consequently $L(2p,\tilde q)$ is rank-$J$ balanced iff $L(p,q)$ is.  The
analogous transfer of $\eta_J=0$ fails: $\eta_3(L(17,3))=0$ but
$\eta_3(L(34,3))=-4$.
\end{lemma}

\begin{proof}
The same $A,B$ of \eqref{eq:AB} define both actions.  Since
$w_L\equiv w_R\equiv n\pmod2$ and $A+B=1$, one has
$Aw_L+Bw_R\equiv n\pmod2$.  For odd $n$ no weight pair is invariant
modulo $2p$; for even $n$ divisibility by $p$ is equivalent to
divisibility by $2p$.  If $L(p,q)$ is balanced so is its double; if the
double is balanced, each residue class modulo the odd number $p$ contains
infinitely many even $n$, and the affine function
\eqref{eq:affine-mod-p} vanishing at them vanishes identically.  The two
eta values quoted are exact computations from \eqref{eq:master}.
\end{proof}

\begin{remark}
For $p=2d$ with $d$ odd satisfying \eqref{eq:clifford}, the rank-$J$
verdict has period $d$ in $J$, sharper than the period $p$ of
Corollary~\ref{cor:criterion-consequences}(b): balance at $2d$ is balance
at the odd parent $d$, whose verdict has period $d$.
\end{remark}

\section{The classification of balanced triples}
\label{sec:classification}

\begin{theorem}[Classification of level-wise balance]
\label{thm:classification}
Let $p\geq3$, $q$ a unit modulo $p$ and $J\geq1$.  Then $L(p,q)$ has
equal numbers of positive- and negative-helicity rank-$J$ TT modes at
every level if and only if one of the following holds:
\begin{enumerate}[label=(\alph*)]
\item \textup{(mirror)} $q^2\equiv-1\pmod p$;
\item \textup{(silver family)} $p\mid4J^2-2$ and $q\equiv\pm(2J\pm1)\pmod p$;
\item \textup{(blind rank)} $p\mid2J$ and the Clifford condition
\eqref{eq:clifford} holds.
\end{enumerate}
\end{theorem}

Sufficiency is Corollary~\ref{cor:criterion-consequences}(c)--(e).  The
proof of necessity occupies the rest of this subsection.  One step of it,
the exhaustion of $105$ finite congruence systems, is certificate-assisted:
the hand reduction below covers the structure, and the complete case list
was verified by integer Gr\"obner-basis certificates whose meaning is made
explicit in Remark~\ref{rem:certificate}.

\begin{proof}[Proof of necessity]
\emph{Step 0 (reductions).}  By Theorem~\ref{thm:criterion}, balance
requires \eqref{eq:clifford}, which fails if $3\mid p$ or $4\mid p$
(Proposition~\ref{prop:existence}).  If $p\equiv2\pmod4$, write $p=2d$
with $d$ odd.  By Lemma~\ref{lem:parity-projection} balance at $2d$ is
equivalent to balance at $d$ with $q\bmod d$; and each of (a), (b), (c) at
$2d$ is equivalent to the same condition at $d$: for (a) because $q$ is
odd, so $q^2\equiv-1\pmod2$ always; for (b) because $4J^2-2$ is even and
$q$, $2J\pm1$ are odd; for (c) because $2d\mid2J$ iff $d\mid J$ iff
$d\mid2J$; and \eqref{eq:clifford} at $2d$ iff at $d$.  It therefore
suffices to treat odd $p\geq5$ with $\gcd(q-1,p)=\gcd(q+1,p)=1$, so that
$2$, $q$ and $q\pm1$ are units modulo $p$.

\emph{Step 1 (matchings).}  Condition \eqref{eq:odd-measure} says that
the multiset of the eight atoms
\begin{gather*}
 a_1=X+q^*,\quad a_2=X-q^*,\quad a_3=Y+q,\quad a_4=Y-q,\\
 a_5=-X-1,\quad a_6=-X+1,\quad a_7=-Y-1,\quad a_8=-Y+1
\end{gather*}
($X=Jq$, $Y=Jq^*$) is invariant under negation in $\Z/p$.  For odd $p$
a nonzero residue is never its own negative, so negation invariance holds
iff the nonzero atoms pair off as $\{a,-a\}$ and the zero atoms are even
in number; that is, iff the eight labelled atoms admit a perfect matching
all four of whose pair sums are $\equiv0\pmod p$.  There are $105$
labelled perfect matchings of eight points.

\emph{Step 2 (pair sums).}  Multiplying each pair sum by the unit $q$
gives an element of $\Z[J,q]$:
\begin{center}
\begin{tabular}{@{}lll@{}}
\toprule
pairs & $q\times$(pair sum) & condition after cancelling units\\
\midrule
$12$, $56$ & $\pm2Jq^2$ & $J\equiv0$\\
$34$, $78$ & $\pm2J$ & $J\equiv0$\\
$13$ & $(J+1)(q^2+1)$ & \\
$24$ & $(J-1)(q^2+1)$ & \\
$14$ & $(J-1)q^2+(J+1)$ & \\
$23$ & $(J+1)q^2+(J-1)$ & \\
$58$, $67$ & $-J(q^2+1)$ & \\
$57$ & $-(J(q^2+1)+2q)$ & \\
$68$ & $-(J(q^2+1)-2q)$ & \\
$17$ & $(q-1)(J(q+1)-1)$ & $J(q+1)\equiv1$\\
$18$ & $(q+1)(J(q-1)+1)$ & $J(q-1)\equiv-1$\\
$27$ & $(q+1)(J(q-1)-1)$ & $J(q-1)\equiv1$\\
$28$ & $(q-1)(J(q+1)+1)$ & $J(q+1)\equiv-1$\\
$35$ & $(q-1)(q-J(q+1))$ & $J(q+1)\equiv q$\\
$36$ & $(q+1)(q-J(q-1))$ & $J(q-1)\equiv q$\\
$45$ & $-(q+1)(J(q-1)+q)$ & $J(q-1)\equiv-q$\\
$46$ & $-(q-1)(J(q+1)+q)$ & $J(q+1)\equiv-q$\\
$15,16,25,26,37,38,47,48$ & $\pm(q-1)$, $\pm(q+1)$, $\pm q(q\pm1)$ & a unit: cannot vanish\\
\bottomrule
\end{tabular}
\end{center}
The eight ``forbidden'' pairs in the last row are units under the
standing hypotheses, so any matching containing one of them is
contradictory.  This leaves $33$ matchings.

\emph{Step 3 (case analysis).}  Call a pair $ij$ a $P$--$N$ pair if
$i\leq4<j$.
\begin{itemize}
\item \emph{Case A: no $P$--$N$ pair ($9$ matchings).}  The pairs are
drawn from $\{12,13,14,23,24,34\}$ and $\{56,57,58,67,68,78\}$.  Each
system reduces, by unit-coefficient linear combinations of the rows of
the table, to $J\equiv0$, to $q^2+1\equiv0$, or to a contradiction of the
form $2q\equiv0$, $4q\equiv0$ or $q^2\equiv1$.
\item \emph{Case B: four $P$--$N$ pairs ($4$ matchings).}  Avoiding the
forbidden pairs, the four cross pairs are one of $17,18,27,28$ paired
with one of $35,36,45,46$, and the complementary pairs.  Each system
contains $J(q\pm1)\equiv1$ together with $J(q\pm1)\equiv-1$ for the same
sign: contradiction since $2$ is a unit.
\item \emph{Case C: two $P$--$N$ pairs ($20$ matchings).}  If the two
cross pairs share a block (subcases C1, C2), the system contains a
contradiction of the same kind as in Case B.  Otherwise (subcase C3, $16$
matchings) one cross pair comes from $\{1,2\}\times\{7,8\}$ and one from
$\{3,4\}\times\{5,6\}$.  Their two conditions $J(q\pm1)\equiv c_1$,
$J(q\mp1)\equiv c_2$ with $c_1,c_2\in\{\pm1,\pm q\}$ are of three kinds.
If both involve the same factor $J(q+1)$ or $J(q-1)$ with different
right sides, they force $q\equiv\pm1$, a contradiction (eight
matchings).  Otherwise subtracting them gives a linear relation
$2J\equiv\pm1\pm q$.  For four matchings, combining the two conditions
then gives $2J^2\equiv1$, which is exactly (b) with $q\equiv\pm(2J\pm1)$.
For the remaining four the linear relation is $2J\equiv1-q$, $1+q$,
$-1-q$ or $-1+q$, which does not by itself make $J$ silver; there the
two remaining same-block pairs are $(J\mp1)K$ and $JK$ with $K=q^2+1$,
so their difference forces $K\equiv0$, that is (a), with the linear
relation as an additional necessary condition on that matching
(Remark~\ref{rem:certificate}).
\end{itemize}
Every step divides only by $2$, $q$ and $q\pm1$, which are units modulo
every odd $p$ satisfying \eqref{eq:clifford}, prime or composite.  Hence
whenever some matching has all four pair sums $\equiv0\pmod p$, one of
$J\equiv0$ (that is, $p\mid J$, so $p\mid2J$, which is (c) since
\eqref{eq:clifford} holds), $q^2\equiv-1$ (that is, (a)), or
$2J^2\equiv1$ with $q\equiv\pm(2J\pm1)$ (that is, (b), because $p$ odd
and $p\mid2J^2-1$ give $p\mid4J^2-2$) holds.
\end{proof}

\begin{remark}[The certificate]
\label{rem:certificate}
The hand analysis of Step~3 is a case list, and it was checked
mechanically as follows.  For each of the $105$ labelled matchings, let
$I\subset\Z[J,q]$ be the ideal generated by its four polynomials
$q\times(\text{pair sum})$, and let $I^{\mathrm{sat}}$ be its saturation by
the unit $2q(q^2-1)$, computed as the elimination ideal of
$I+\bigl(t\cdot2q(q^2-1)-1\bigr)$ in $\Z[J,q,t]$ by a Gr\"obner basis with
integer coefficient domain.  A matching is \emph{certified} contradictory
if $I^{\mathrm{sat}}=(1)$, and certified of class (a), (b) or (c) if the
generators $q^2+1$, resp.\ $\{2J^2-1,\ q-s\}$ ($s$ the named root
$\pm(2J\pm1)$), resp.\ $J$ reduce to $0$ modulo $I^{\mathrm{sat}}$ over the
integers.  Membership of $f$ in $I^{\mathrm{sat}}$ means that
$(2q(q^2-1))^Nf\in I$ for some $N$, an identity in $\Z[J,q]$; reducing
modulo any odd $p$ satisfying \eqref{eq:clifford} then gives
$f\equiv0\pmod p$ whenever the four pair sums vanish.  The certificate
is therefore a set of \emph{implications} (vanishing pair sums force the
class generators to vanish), tested matching by matching; it is not a
list of $105$ exported ideal bases, and one representative basis per
class (the first of its class) was retained.  The certified classes
are: $93$ contradictions (representative basis $\{1\}$); two matchings
of class (c), $\{12,34,56,78\}$ and $\{12,34,58,67\}$ (representative
$\{J\}$); one matching, $\{13,24,56,78\}$, of both classes (a) and (c)
(representative $\{1+q^2,J\}$: the ideal forces both conditions, not
their union); five of class (a), $\{13,24,58,67\}$, $\{17,24,36,58\}$,
$\{18,24,35,67\}$, $\{13,27,46,58\}$, $\{13,28,45,67\}$ (representative
$\{1+q^2\}$, the basis of the first of them); and four of class (b), one
for each root, with the bases
\begin{center}
\begin{tabular}{@{}lll@{}}
\toprule
matching & root & certified Gr\"obner basis\\
\midrule
$\{17,23,45,68\}$ & $q\equiv2J-1$ & $\{q^2+2q-1,\ 2J-q-1,\ J(q+1)-1\}$\\
$\{18,23,46,57\}$ & $q\equiv-(2J-1)$ & $\{q^2-2q-1,\ 2J+q-1,\ J(q+1)+q\}$\\
$\{14,27,35,68\}$ & $q\equiv2J+1$ & $\{q^2-2q-1,\ 2J-q+1,\ J(q+1)-q\}$\\
$\{14,28,36,57\}$ & $q\equiv-(2J+1)$ & $\{q^2+2q-1,\ 2J+q+1,\ J(q+1)+1\}$\\
\bottomrule
\end{tabular}
\end{center}
Each (b) basis says exactly ``$q$ is the named root and $2J^2\equiv1$''.
The representative basis $\{1+q^2\}$ is the saturated ideal of
$\{13,24,58,67\}$ only.  For the four class-(a) matchings with two cross
pairs the ideal also contains a linear relation between $q$ and $J$:
\begin{center}
\begin{tabular}{@{}ll@{}}
\toprule
matching & relations forced, besides $q^2+1\equiv0$\\
\midrule
$\{17,24,36,58\}$ & $2J+q-1\equiv0$\\
$\{18,24,35,67\}$ & $2J-q-1\equiv0$\\
$\{13,27,46,58\}$ & $2J+q+1\equiv0$\\
$\{13,28,45,67\}$ & $2J-q+1\equiv0$\\
\bottomrule
\end{tabular}
\end{center}
(for the first, the pairs $17$ and $36$ give $J(q+1)\equiv1$ and
$J(q-1)\equiv q$, whence $2J\equiv1-q$; the pairs $24$ and $58$ give
$(J-1)(q^2+1)\equiv J(q^2+1)\equiv0$, whence $q^2+1\equiv0$).  That the
one-generator basis is not the full ideal of these four is witnessed at
$p=5$, $q=2$, $J=5$: there $q^2+1\equiv0$ and $2q(q^2-1)$ is a unit, yet
the pair sums of $\{17,24,36,58\}$ are $(2,0,3,0)$ modulo $5$.  This
concerns only the description of the certificate; the implication to
class (a), which is all that the classification uses, holds for all
five.  The semantic check of Appendix~\ref{app:verification} confirms on
concrete moduli that the criterion is equivalent to
``\eqref{eq:clifford} and some matching with vanishing pair sums'', that
every such matching is one of the twelve above, and that the criterion
is equivalent to (a) or (b) or (c).

\emph{Elementary hand check.}  Exhaust the $33$
matchings without a forbidden pair by hand, dividing only by the units
$2$, $q$, $q\pm1$ and never by $J$ or by $q^2+1$: the $9$ matchings
without a cross pair (the positive pairings $(12,34)$, $(13,24)$,
$(14,23)$ against the negative pairings $(56,78)$, $(57,68)$, $(58,67)$;
the first forces $J\equiv0$ and is compatible with $(56,78)$ and
$(58,67)$ but not with $(57,68)$; the second forces $q^2+1\equiv0$ by
subtraction and is compatible with $(56,78)$ only together with
$J\equiv0$, with $(58,67)$ without restriction, and not with $(57,68)$;
the third forces $2(q^2-1)\equiv0$, impossible), giving $4$ possible
matchings and $5$ contradictions; the $4$ matchings with four cross
pairs and the $4$ with two cross pairs in the same pair of blocks, all
contradictions (equal factors $J(q\pm1)$ with opposite right sides); and
the $16$ mixed-block matchings, of which $8$ force $q\equiv\pm1$, $4$ are
the mirror matchings of the table above, and $4$ are the silver
matchings (cross pairs $(17,45)$, $(18,46)$, $(27,35)$, $(28,36)$, giving
$q\equiv2J-1$, $1-2J$, $2J+1$, $-2J-1$ with $2J^2\equiv1$).  In all, the
$12$ possible matchings and $21$ contradictions of the certified list,
which with the $72$ matchings containing a forbidden pair exhaust the
$105$.  This gives the
elementary proof of the case list; the certificate remains as
mechanical corroboration.
\end{remark}

\begin{corollary}
\label{cor:classification-consequences}
Let $p\geq3$.
\begin{enumerate}[label=(\alph*)]
\item A chiral lens space is rank-$J$ balanced only through the blind
ranks (c) or the silver family (b).  For odd $p$, (b) allows at most four
residue classes of $J$ modulo $p$ (namely $2J\equiv\pm q\pm1$); for
$p\equiv2\pmod4$ the same holds modulo $p/2$.
\item If $p$ is odd and $L(p,q)$ is chiral and rank-$J$ balanced, then
$p\mid J$ or $\gcd(J,p)=1$.
\end{enumerate}
\end{corollary}

\begin{proof}
(a) is immediate from the theorem.  (b) In case (c), $p\mid2J$ and $p$ odd
give $p\mid J$; in case (b), $p\mid2J^2-1$ and $\gcd(J,2J^2-1)=1$ give
$\gcd(J,p)=1$.
\end{proof}

\begin{corollary}[Rank one]
\label{cor:rank-one}
For every $p\geq1$ and every unit $q$, the rank-one (coexact curl)
spectrum of $L(p,q)$ is balanced at every level if and only if
$q^2\equiv-1\pmod p$, i.e.\ iff $L(p,q)$ is amphicheiral.
\end{corollary}

\begin{proof}
For $p\geq3$ this is Theorem~\ref{thm:classification} at $J=1$: (b)
would need $p\mid2$ and (c) would need $p\mid2$.  For $p\leq2$ both sides
hold.  A direct proof valid for all $p$: if $q^2\equiv-1$ the
orientation-reversing isometry $T$ of Definition~\ref{def:terms}
exchanges the two branches.  Conversely balance forces
$\eta_1(L(p,q))=-4s(q,p)=0$, and $s(q,p)=0$ iff $q^2\equiv-1\pmod p$
\cite[p.~36]{RademacherGrosswald1972}.  For completeness: if $s(q,p)=0$
then Dedekind reciprocity multiplied by $12pq$ gives
$12pq\,s(p,q)=p^2+q^2+1-3pq$, whose left side is a multiple of $p$ because
$6q\,s(p,q)\in\Z$, so $q^2+1\equiv0\pmod p$; if $q^*\equiv-q$ then
$s(q,p)=s(q^*,p)=s(-q,p)=-s(q,p)$.
\end{proof}

This rank-one statement is classical by composition: the eta invariant
of coexact curl is the signature eta (Millson, as recalled in
\cite{Bar2019}), which
on $L(p,q)$ is $-4s(q,p)$ \cite{APS1975,Katase1987,Jeffrey1992}; the
Dedekind sum vanishes iff $q^2\equiv-1$
\cite{RademacherGrosswald1972,Meyer2005,Girstmair2016}; and
$q^2\equiv-1$ is the amphicheirality criterion
\cite{PrzytyckiYasuhara2003,Boldt2017}.  The composed sentence itself we
did not find in print; its nearest published relative is
\cite[Rem.~4.13]{Boldt2017} for the Dirac operator, proved there for
prime $p$ and conjectured for odd $p$.  The rank-$J\geq2$ statements of
Theorems~\ref{thm:criterion} and \ref{thm:classification}, and the
families (b) and (c), were not found in the literature we searched; the
helicity-\emph{summed} spin-$s$ multiplicities on $L(p,q)$ as sums over
$\Z_p$ of products of $SU(2)$ characters are standard
\cite{CastroLashkariMaloney2011Farey,Basu2015,MartinSvesko2020}.  This
is a bounded statement about a search, not a claim of absence.

\begin{example}[Zero eta without balance]
\label{ex:zero-not-balanced}
The four zero families of \cite[Thm.~7.1]{MahmoodP1Collection2026} at
$p=2J^2-1$ are $q=J$, $J+1$, $2J-1$ and $J(J-1)$ and their orbit images.
The parametrization $q=2J-1$ (with its inverse $2J+1$ and their
negatives) is uniformly balanced, at every $J$; the other three are not
balanced in general, but their members can overlap the balanced or
mirror loci, and the named families are not disjoint.  For $J=3$,
$p=17$: $\eta_3(L(17,q))=0$ for $q\in\{3,6,11,14\}$ (the orbit of
$q=J=3$), but $L(17,3)$ has $(N^+_3,N^-_3)=(1,3)$ at level $n=2$;
$L(17,4)$ (the orbit of $q=J+1$) is balanced, but only because
$4^2\equiv-1\pmod{17}$, i.e.\ it is mirror; and $L(17,5)$ (the orbit of
$q=2J-1$) is balanced and chiral.  At $J=2$, $p=7$, the member $q=J=2$
of the first family lies in the balanced silver orbit $\{2,3,4,5\}$
(Proposition~\ref{prop:C3-balance}) and is balanced too, whereas
$L(17,3)$ above is a genuine regularization-only zero.  Similarly
$\eta_2(L(35,q))=0$ for
$q\in\{8,13,22,27\}$ with $(N^+_2,N^-_2)=(2,0)$ at $n=3$ on $L(35,8)$.
Balance is rank-specific: $L(17,5)$ is balanced at $J=3$ but has
$(N^+_2,N^-_2)=(2,0)$ at $n=1$.
\end{example}

\begin{example}[Small balanced chiral spaces]
\label{ex:balanced-table}
For $p\leq260$ and $J\leq8$, the chiral rank-$J$ balanced spaces are
exactly the following (exact computation; the rows $J\geq5$ were
predicted from Theorem~\ref{thm:classification}'s families before being
computed, see Appendix~\ref{app:verification}):
\begin{center}
\begin{tabular}{@{}ll@{}}
\toprule
$J$ & balanced chiral $(p:\ q)$, $p\leq260$\\
\midrule
$2$ & $7{:}\{2,3,4,5\}$; $14{:}\{3,5,9,11\}$\\
$3$ & $17{:}\{5,7,10,12\}$; $34{:}\{5,7,27,29\}$\\
$4$ & $31{:}\{7,9,22,24\}$; $62{:}\{7,9,53,55\}$\\
$5$ & $7{:}\{2,3,4,5\}$; $14{:}\{3,5,9,11\}$; $49{:}\{9,11,38,40\}$; $98{:}\{9,11,87,89\}$\\
$6$ & $71{:}\{11,13,58,60\}$; $142{:}\{11,13,129,131\}$\\
$7$ & $7{:}\{2,3,4,5\}$; $14{:}\{3,5,9,11\}$ (blind rank); $97{:}\{13,15,82,84\}$; $194{:}\{13,15,179,181\}$\\
$8$ & $127{:}\{15,17,110,112\}$; $254{:}\{15,17,237,239\}$\\
\bottomrule
\end{tabular}
\end{center}
At $J=5$ the divisor $7$ of $4\cdot25-2=98$ enters family (b); at $J=7$
the spaces $L(7,q)$, $L(14,q)$ enter through the blind rank $7\mid2J$.
Rank scans at fixed space: $L(7,3)$ and $L(14,3)$ are balanced exactly
at $J\in\{2,5,7,9,12,14,16\}$ for $J\leq16$; $L(17,5)$ at
$\{3,14,17,20,31,34\}$ for $J\leq36$; $L(11,2)$ at $\{11,22\}$ for
$J\leq24$; $L(9,2)$ at no rank (it fails \eqref{eq:clifford});
$L(23,4)$ at $\{9,14,23\}$ for $J\leq25$; $L(31,7)$ at $\{4,27,31\}$ for
$J\leq31$;
$L(49,9)$ at $\{5,44,49\}$ for $J\leq50$.
\end{example}

\section{Divisor decomposition of the eta invariant}
\label{sec:decomposition}

Level-wise balance is a strictly smaller phenomenon than $\eta_J=0$
(Example~\ref{ex:zero-not-balanced}).  To study the zeros of $\eta_J$
themselves we first localize $p\,\eta_J$ over the divisors of $p$.

\subsection{Per-element regularized sums}

For a nonidentity element $g^k$ of $L(p,q)$ let
\begin{equation}
 \eta_k(p,q,J)=\sum_{n\geq0}^{\text{reg}}E_k(n)\,\nu^{-s}\Big|_{s=0},
 \qquad\nu=n+J+1,
 \label{eq:per-element-eta}
\end{equation}
the zeta-regularized level sum of the signed character
\eqref{eq:Ek}.  Since $E_k(n)$ is, for fixed $k$, a finite linear
combination of $\nu^j\omega^\nu$ with $j\leq1$ and $\omega$ a $2p$th root
of unity, its Dirichlet series is a finite combination of Hurwitz zeta
functions, is regular at $s=0$, and equals there the Abel value of the
series \cite[Lem.~5.4]{MahmoodP1Collection2026}; the identity element
contributes zero at every level, and $\omega=1$ never occurs for
$k\neq0$.  Hence
\begin{equation}
 p\,\eta_J(L(p,q))=\sum_{k=1}^{p-1}\eta_k(p,q,J).
 \label{eq:eta-as-element-sum}
\end{equation}
Writing $\theta=\theta_k$, $N=J+1$, and
\[
 S_{\sin}(a)=\tfrac12\cot\tfrac a2-\sum_{\nu=1}^{N-1}\sin(\nu a),\qquad
 S_{n\cos}(a)=-\frac1{4\sin^2(a/2)}-\sum_{\nu=1}^{N-1}\nu\cos(\nu a)
\]
for the Abel values of $\sum_{\nu\geq N}\sin(\nu a)$ and
$\sum_{\nu\geq N}\nu\cos(\nu a)$ ($a\notin2\pi\Z$), the proof of
Theorem~\ref{thm:kernel} gives the explicit terms
\begin{align}
 \eta_k&=\frac{2\bigl[S_{\sin}(\theta)\sin(Jq\theta)-S_{\sin}(q\theta)\sin(J\theta)\bigr]}
 {\cos(q\theta)-\cos\theta}
 &&(k\in\mathcal N),\label{eq:generic-element}\\
 \eta_k&=\pm\frac{2\bigl[S_{n\cos}(\theta)\sin(J\theta)-J\,S_{\sin}(\theta)\cos(J\theta)\bigr]}
 {\sin\theta}
 &&(k\in D^{\pm*}),\label{eq:clifford-element}\\
 \eta_{p/2}&=0&&(p\text{ even}).\nonumber
\end{align}

\begin{lemma}[Trace form for nonsingular $q$]
\label{lem:trace-form}
Let $p\geq3$ be odd and let the unit $q$ satisfy
$\gcd(q-1,p)=\gcd(q+1,p)=1$ (so every
nonidentity element is generic).  Then, with $\varphi_k=\pi k/p$,
\begin{equation}
 p\,\eta_J(L(p,q))=\sum_{k=1}^{p-1}
 \frac{\cos((2J+1)\varphi_k)\dfrac{\sin(2Jq\varphi_k)}{\sin\varphi_k}
 -\cos((2J+1)q\varphi_k)\dfrac{\sin(2J\varphi_k)}{\sin(q\varphi_k)}}
 {\cos(2q\varphi_k)-\cos(2\varphi_k)},
 \label{eq:E1}
\end{equation}
and, after the reindexing $k\mapsto q^*k$ in the second term,
\begin{equation}
 p\,\eta_J(L(p,q))=\sum_{k=1}^{p-1}
 \cK_J(\theta_k)\,\frac{\cos\bigl((J+\tfrac12)\theta_k\bigr)}{\sin(\theta_k/2)}.
 \label{eq:trace-form}
\end{equation}
\end{lemma}

\begin{proof}
Sum the formal character of
\cite[Sec.~5.1]{MahmoodP1Collection2026},
\[
 D^{\mathrm{form}}_J(m;\ell)=\frac{\sin(2m\beta)\sin(2J\alpha)-\sin(2m\alpha)\sin(2J\beta)}
 {\sin(\alpha+\beta)\sin(\beta-\alpha)},\qquad
 \alpha=\frac{\pi\ell q}p,\quad\beta=\frac{\pi\ell}p,
\]
over the physical modes $m\geq J+1$ using the Abel
value $\sum_{m\geq J+1}\sin(2mx)=\cos((2J+1)x)/(2\sin x)$ and
$\sin(\alpha+\beta)\sin(\beta-\alpha)=\tfrac12(\cos2\alpha-\cos2\beta)$.
Summing over $m\geq J+1$ directly removes the sub-threshold modes, so no
$L_J$ term appears.  With $\varphi=\beta$ and $\alpha=q\varphi$ this is
\eqref{eq:E1}, which is \eqref{eq:generic-element} summed over $k$.  The
reindexing is Step~3 of Theorem~\ref{thm:kernel}; in terms of
$x=\e^{i\theta}$ the summand of \eqref{eq:trace-form} is
$f(x)=C(x)[A_q(x)+A_{q^*}(x)]$ with
\begin{equation}
 C(x)=\frac{x^{J+1}+x^{-J}}{x-1},\qquad
 A_r(x)=\frac{x^{Jr}-x^{-Jr}}{x^r+x^{-r}-x-x^{-1}},
 \label{eq:C-A}
\end{equation}
which is even under $x\mapsto1/x$.
\end{proof}

\subsection{Primitive parts}

\begin{theorem}[Divisor decomposition]
\label{thm:decomposition}
For $p\geq2$, every unit $q$ modulo $p$ and $J\geq1$,
\begin{equation}
 p\,\eta_J(L(p,q))=\sum_{\substack{d\mid p\\ d>1}}P_d(q\bmod d,\,J),
 \qquad
 P_d(q,J)=\sum_{k\in(\Z/d)^\times}\eta_k(d,q,J),
 \label{eq:decomposition}
\end{equation}
where $P_d$, the \emph{primitive part} of order $d$, is the regularized
contribution of the elements of exact order $d$ and depends only on
$(d,\,q\bmod d,\,J)$.  M\"obius inversion recovers the primitive parts
from the eta invariants of the divisor lens spaces:
\begin{equation}
 P_d(q,J)=\sum_{\substack{e\mid d\\ e>1}}\mu(d/e)\,e\,\eta_J(L(e,q\bmod e)).
 \label{eq:moebius}
\end{equation}
\end{theorem}

\begin{proof}
Group the nonidentity elements in \eqref{eq:eta-as-element-sum} by their
exact order $d$.  An element of order $d$ in $L(p,q)$, $k=(p/d)k'$ with
$k'\in(\Z/d)^\times$, acts on $S^3$ as the rotation
$(\e^{2\pi ik'/d},\e^{2\pi iqk'/d})$, which is the element $g^{k'}$ of
$L(d,q\bmod d)$; its signed character at every level, hence its
regularized sum, depends only on $(d,q\bmod d,J,k')$ (the character
products are lift-independent by \cite[Lem.~5.2]{MahmoodP1Collection2026}).
The finite sum over group elements commutes with the continuation, and
the identity contributes nothing ($H_1=0$).  Applying
\eqref{eq:decomposition} to each divisor $e$ and inverting gives
\eqref{eq:moebius}.
\end{proof}

\begin{corollary}
\label{cor:decomposition}
\begin{enumerate}[label=(\alph*)]
\item At a prime $p$, $p\,\eta_J(L(p,q))=P_p(q,J)$ is a single primitive
part.
\item If $L(p,q)$ is rank-$J$ balanced or mirror, then
$P_d(q\bmod d,J)=0$ for every $d\mid p$.
\item If every primitive part vanishes then $\eta_J=0$.  The converse
fails at composite $p$: for $(p,q,J)=(42,11,9)$ one has $\eta_9=0$ with
$D_9(0)=2$, and the primitive parts are $P_3=12$, $P_6=36$, $P_{21}=-54$,
$P_{42}=6$, $P_2=P_7=P_{14}=0$.
\end{enumerate}
\end{corollary}

\begin{proof}
(a) is the case of a single nontrivial divisor.  (b) Balance and the
mirror condition descend to every divisor
(Corollary~\ref{cor:criterion-consequences}(a); $q^2\equiv-1\pmod d$),
where $\eta_J=0$, so every M\"obius sum \eqref{eq:moebius} vanishes.
(c) The first claim is \eqref{eq:decomposition}; the example is an exact
computation from \eqref{eq:master} and \eqref{eq:moebius}.
\end{proof}

\begin{proposition}[Shift invariance and the first-moment form]
\label{prop:first-moment}
Fix $p,q,J$ and write $D_J(r+pk)=v_r+ku_r$ as in \eqref{eq:affine-mod-p}.
For any $c>0$ the value at $s=0$ of $\sum_{n\geq0}D_J(n)(n+c)^{-s}$ is
\[
 \sum_{r=0}^{p-1}\Bigl[v_r\Bigl(\tfrac12-a_r\Bigr)
 +u_r\Bigl(\tfrac{a_r^2}2-\tfrac1{12}\Bigr)\Bigr],\qquad a_r=\frac{r+c}p,
\]
and it is independent of $c$.  In particular, under
\eqref{eq:clifford} (where $u_r=0$ and $D_J$ is $p$-periodic with
$\sum_rv_r=0$),
\begin{equation}
 \eta_J(L(p,q))=-\frac1p\sum_{r=0}^{p-1}r\,D_J(r).
 \label{eq:first-moment}
\end{equation}
\end{proposition}

\begin{proof}
The Hurwitz evaluation is \cite[Prop.~4.1]{MahmoodP1Collection2026} with
period $p$.  By \eqref{eq:FG}, $pD_J(n)=\nu F(\nu)+G(\nu)$ where $F$ and
$G$ are finite Fourier sums over nonidentity elements, so both have mean
zero over a period: $\sum_ru_r=0$ and $p\sum_rv_r=\sum_r(r+J+1)u_r$.
Shifting $c\mapsto c+h$ changes the value by
$\frac h{p^2}\bigl[-p\sum_rv_r+\sum_r(r+c)u_r\bigr]+\frac{h^2}{2p^2}\sum_ru_r$,
which vanishes by these two identities.  Under \eqref{eq:clifford},
$F=0$, so $u_r=0$, and $\sum_rv_r=0$ gives \eqref{eq:first-moment}.
\end{proof}

Thus at a prime $p$ with $q\not\equiv\pm1$, $\eta_J=0$ requires only a
vanishing first moment of the periodic sequence $D_J(0),\ldots,D_J(p-1)$,
whereas balance requires every term to vanish.  This is the algebraic
reason why the classification of Section~\ref{sec:classification} does
not classify the zeros of $\eta_J$.

\subsection{Observations at prime moduli}

The following facts were found by exact computation and are recorded as
observations; the theorems that explain them are in
Sections~\ref{sec:congruences} and \ref{sec:silver}.

\begin{observation}
\label{obs:prime-zeros}
Over all primes $41\leq p\leq199$, all units $q$ and $1\leq J\leq12$
($47{,}952$ rows), and all primes $41\leq p\leq101$, all $q$ and
$13\leq J\leq(p-1)/2$ ($23{,}278$ rows), $\eta_J(L(p,q))$ has $1{,}004$
zeros.  Every zero is one of: a mirror space; a member of the silver family
(b); a member of the four families of \cite[Thm.~7.1]{MahmoodP1Collection2026}
at $p=2J'^2-1$ with $J\equiv\pm J'$; or one of $196$ further zeros in
$49$ orbits under $q\mapsto\pm q^{\pm1}$.  Every one of these $49$ orbits
lies at a silver rank $2J^2\equiv1\pmod p$; the nine silver pairs
$(p,J)$ in range are $(41,12)$, $(71,6)$, $(97,7)$, $(127,8)$, $(199,10)$,
$(47,20)$, $(73,16)$, $(79,35)$, $(89,32)$, and every one of them carries
such zeros.  Among them, $q\equiv J$ is a zero at every silver pair in
range (the first family of \cite{MahmoodP1Collection2026}, at primes
dividing $2J^2-1$ rather than equal to it), whereas the analogous
extensions of $q=J+1$ and $q=J(J-1)$ fail; at $(97,7)$ and $(199,10)$ the
cube roots of unity modulo $p$ are also zeros.  At silver ranks every
value in range is an even integer; at other ranks $\eta_J$ is in general
a non-integral rational.  For example $\eta_{12}(L(41,12))=0$ with
$D_{12}(3)=-2$, while $\eta_{12}(L(41,1))=-76$.
\end{observation}

\begin{observation}
\label{obs:silver-blind}
On primes never inspected before the following three statements were
formulated: for every prime $211\leq p\leq331$, every $q$ and $J\leq12$
off the silver ranks ($67{,}032$ rows) the $238$ zeros are exactly the
mirror rows; for every non-silver rank $13\leq J\leq(p-1)/2$ at
$p\in\{211,223\}$ ($41{,}286$ rows) there is no zero, as both primes are
$\equiv3\pmod4$; at the silver rank of each of the $31$ primes
$p\equiv\pm1\pmod8$ in $[211,599]$ and every $q$ ($12{,}210$ rows) every
value is an even integer, and $q=J$ is a zero at all $31$ primes.  At the
adjacent rank $J+1$, the control computation on the first eight of these
primes ($223\leq p\leq281$) has $1{,}992$ nonintegral values among
$2{,}000$ rows; only the eight mirror rows are integral, as
Corollary~\ref{cor:E}(a) explains.  At the silver rank $J$, the residue
$q=J+1$ fails to be a zero at $22$ of the $31$ primes.  Descriptively, at
these silver ranks $471$ of $3{,}076$ $q$-orbits ($15.3\%$) are zeros; the
zero fraction is highest for small silver ranks ($33\%$ at $p=241$,
$J=11$; $27\%$ at $p=449$, $J=15$; $25\%$ at $p=337$, $J=13$), and the
extreme values occur at $q=\pm1$ (up to $|\eta|=15{,}640$ at $p=569$).
\end{observation}

\begin{observation}[Composite moduli]
\label{obs:composite}
Over composites $41\leq p\leq120$, all units $q$ and $J\leq10$
($25{,}180$ rows), the zeros of $\eta_J$ fall into the classes: balanced
(Theorem~\ref{thm:classification}) $196$; the families of
\cite[Thm.~7.1]{MahmoodP1Collection2026} $12$; all primitive parts zero
but neither of the above $24$; and cross-cancellations of $2$, $3$, $4$
nonzero primitive parts, $44$, $16$, $28$ respectively.
\end{observation}

\section{Congruences for the eta invariant}
\label{sec:congruences}

\subsection{Two elementary inputs valid at every modulus}

\begin{lemma}[Low-mode parity]
\label{lem:low-mode-parity}
For every $p\geq2$, every unit $q$ and all $a,b\geq1$ with
$a\equiv b\pmod2$,
$I_{a,b}(p,q)\equiv\ind_{a,b\text{ both odd}}\pmod2$.  Hence
$I_{a,b}-I_{b,a}$ is even and $L_J(p,q)$ is even.
\end{lemma}

\begin{proof}
The involution $(r,t)\mapsto(-r,-t)$ of $W_{a-1}\times W_{b-1}$ preserves
the congruence $(1+q)r+(1-q)t\equiv0\pmod{2p}$.  Its only fixed point is
$(0,0)$, which lies in $W_{a-1}\times W_{b-1}$ exactly when $a$ and $b$
are both odd.
\end{proof}

\begin{lemma}[Dedekind-sum facts at every modulus]
\label{lem:dedekind-facts}
Let $p\geq2$ and $1\leq q<p$ with $\gcd(q,p)=1$.  For $1\leq k<p$ put
$r_k=\langle qk\rangle_p\in\{1,\ldots,p-1\}$, $n_k=\lfloor qk/p\rfloor$,
and
\[
 R=\sum_{k=1}^{p-1}k\,r_k,\qquad
 S_2=\sum_{k=1}^{p-1}k^2=\frac{p(p-1)(2p-1)}6,\qquad
 \Lambda=\sum_{k=1}^{p-1}n_k^2 .
\]
Then
\begin{align}
 2qR&=(q^2+1)S_2-p^2\Lambda,\label{eq:square-identity}\\
 s(q,p)&=\frac R{p^2}-\frac{p-1}4,\qquad
 T:=12p\,s(q,p)=\frac{12R}p-3p(p-1)\in2\Z,\label{eq:T}\\
 qT&\equiv q^2+1\pmod{\gcd(3,p)\,p},\qquad
 T\equiv q+q^*\pmod p.\label{eq:rademacher-congruence}
\end{align}
If $p\geq5$ is prime, moreover $p\mid R$ and $2p\,s(q,p)\in\Z$.
\end{lemma}

\begin{proof}
The residues $r_k$ permute $\{1,\ldots,p-1\}$.  Expanding
$r_k=qk-pn_k$ in $\sum r_k^2=S_2$ and using $\sum k\,n_k=(qS_2-R)/p$
gives \eqref{eq:square-identity}.  Expanding the sawtooth definition of
$s(q,p)$ gives the first part of \eqref{eq:T}.  Multiplying
\eqref{eq:square-identity} by $3$ shows $p\mid6qR$, since
$(q^2+1)(p-1)(2p-1)/2$ is an integer ($p-1$ is even for odd $p$, and
$q^2+1$ is even for even $p$); as $\gcd(q,p)=1$, $p\mid6R$, so $T$ is an
integer, and it is even because $12R/p=2\cdot(6R/p)$ and $3p(p-1)$ are
even.  Multiplying \eqref{eq:square-identity} by $6$, dividing by $p$
and substituting \eqref{eq:T} gives the exact identity
\begin{equation}
 qT-(q^2+1)=p\bigl[(q^2+1)(2p-3)-6\Lambda-3q(p-1)\bigr],
 \label{eq:qT-identity}
\end{equation}
whose bracket is divisible by $3$ when $3\mid p$; this is
\eqref{eq:rademacher-congruence}, the second congruence following on
multiplication by $q^*$.  If $p\geq5$ is prime, $6$ is invertible modulo
$p$ and $R\equiv qS_2\equiv0\pmod p$ (a permutation of the residues
changes nothing modulo $p$), whence $2ps=2R/p-p(p-1)/2\in\Z$.
\end{proof}

The congruence $12pq\,s(q,p)\equiv q^2+1\pmod{\gcd(3,p)p}$ is
Rademacher's \cite[Ch.~3]{RademacherGrosswald1972}; the finite-sum proof
above is included so that nothing below rests on an unverified citation.

\subsection{Prime moduli}

\begin{theorem}[Prime-modulus congruence]
\label{thm:E}
Let $p\geq5$ be prime, $q$ a unit modulo $p$ and $J\geq1$.  Then
\begin{enumerate}[label=(\roman*)]
\item $p\,\eta_J(L(p,q))\in2\Z$;
\item $p\,\eta_J(L(p,q))\equiv-\dfrac{J(2J^2-1)(q+q^{-1})}3\pmod p$.
\end{enumerate}
Equivalently, $\eta_J\equiv J(2J^2-1)\,\eta_1$ modulo $\Z_{(p)}$, where
$\eta_1=-4s(q,p)$ and
$p\eta_1=-T/3\equiv-(q+q^{-1})/3\pmod p$ by
Lemma~\ref{lem:dedekind-facts}.
\end{theorem}

\begin{proof}
Multiply \eqref{eq:master} by $p$.  The Dedekind term becomes
$-2J\cdot(2ps)$, an even integer by Lemma~\ref{lem:dedekind-facts}; the
residue term is $2\sum_a(J-a)[\rho_p(aq)+\rho_p(aq^*)]$, visibly even;
the aliasing term is $2p$ times an integer; and $-pL_J$ is even by
Lemma~\ref{lem:low-mode-parity}.  This is (i).  Modulo $p$ the last two
terms vanish.  Since $\rho_p(k)\equiv-2k\pmod p$ in both cases of its
definition, and $\sum_{a=1}^{J-1}a(J-a)=J(J^2-1)/6$,
\[
 2\sum_{a=1}^{J-1}(J-a)\bigl[\rho_p(aq)+\rho_p(aq^*)\bigr]
 \equiv-4(q+q^*)\sum_{a=1}^{J-1}a(J-a)
 =-\tfrac23J(J^2-1)(q+q^*),
\]
while $-4Jps=-\tfrac J3T\equiv-\tfrac J3(q+q^*)$ by
\eqref{eq:rademacher-congruence}; here $3$ is invertible because
$p\geq5$.  Adding, $p\,\eta_J\equiv-\tfrac J3(q+q^*)[1+2(J^2-1)]
=-J(2J^2-1)(q+q^*)/3$, which is (ii).
\end{proof}

\begin{corollary}[Integrality, denominators, and the zero set off the silver ranks]
\label{cor:E}
Let $p\geq5$ be prime, $q$ a unit and $J\geq1$.
\begin{enumerate}[label=(\alph*)]
\item $\eta_J(L(p,q))\in\Z$ iff $p\mid J(2J^2-1)(q^2+1)$, and then
$\eta_J\in2\Z$.  Otherwise the reduced denominator of $\eta_J$ is exactly
$p$; in particular $\eta_J\neq0$.
\item \textup{(Silver integrality.)}  If $2J^2\equiv1\pmod p$ then
$\eta_J(L(p,q))\in2\Z$ for every unit $q$.
\item \textup{(Dichotomy off the silver ranks.)}  If $2J^2\not\equiv1\pmod p$
and $p\nmid J$, then $\eta_J(L(p,q))=0$ iff $q^2\equiv-1\pmod p$.
\item \textup{(Blind ranks.)}  If $p\mid J$ then $\eta_J(L(p,q))=0$ for
every $q\not\equiv\pm1$, while for $J=np$, $n\geq1$,
\begin{equation}
 \eta_{np}(L(p,\pm1))=\mp\frac{n(p^2-1)}6\neq0,
 \label{eq:blind-clifford}
\end{equation}
and $\eta_{J+p}(L(p,\pm1))-\eta_J(L(p,\pm1))=\mp(p^2-1)/6$ for every
$J\geq1$.
\end{enumerate}
Consequently, for a prime $p\geq5$ and a rank $J$ with
$2J^2\not\equiv1\pmod p$, the zero set of $\eta_J(L(p,\cdot))$ is
completely known: the mirror spaces if $p\nmid J$, and all
$q\not\equiv\pm1$ if $p\mid J$.  At $J=1$, (c) is the classical Dedekind
zero criterion at prime moduli.
\end{corollary}

\begin{proof}
(a) Write $p\,\eta_J=2\Theta$ with $\Theta\in\Z$ by (i).  Then
$\eta_J\in\Z$ iff $p\mid2\Theta$ iff $p\mid\Theta$ ($p$ odd), and then
$\eta_J=2\Theta/p\in2\Z$; otherwise the denominator is exactly $p$.  By
(ii), $p\mid\Theta$ iff $p\mid J(2J^2-1)(q+q^*)$ iff
$p\mid J(2J^2-1)(q^2+1)$ (multiply by the unit $q$).  (b) is (a) with
$p\mid2J^2-1$.  (c) Under the hypotheses $p\nmid J(2J^2-1)$, so by (a)
$\eta_J$ is nonzero unless $p\mid q^2+1$; conversely mirror spaces have
$\eta_J=0$ by \eqref{eq:lens-symmetries}.  (d) The vanishing for
$q\not\equiv\pm1$ is the rank reflection $\eta_p=0$ and periodicity of
\cite[Thm.~6.2]{MahmoodP1Collection2026}, whose hypothesis
$\gcd(1\pm q,p)=1$ holds at a prime for $q\not\equiv\pm1$.  For $q=1$
every element is a Clifford translation and, by
\eqref{eq:clifford-element} and \eqref{eq:C-A}, in the variable
$x=\e^{i\theta_k}$,
\[
 \eta_k=f_{J,1}(x)
 :=\frac{J\,C(x)\,(x^J+x^{-J})-C'(x)\,(x^J-x^{-J})}{x-x^{-1}},
 \qquad C'(x)=x\frac{dC}{dx},
\]
so that $p\,\eta_J(L(p,1))=\Tr_{\Q(\zeta_p)/\Q}f_{J,1}(\zeta_p)$.  At
$J=np$ and a primitive $p$th root of unity, $x^J=x^{-J}=1$ and
$x^{J+1}=x$, so
\[
 f_{J,1}(x)=\frac{2J\,C(x)}{x-x^{-1}}
 =\frac{2J(x+1)}{(x-1)(x-x^{-1})}=\frac{2Jx}{(x-1)^2}
 =2J\bigl(\lambda^{-2}-\lambda^{-1}\bigr),\qquad\lambda=1-x .
\]
The elementary traces $\Tr\lambda^{-1}=(p-1)/2$ and
$\Tr\lambda^{-2}=-(p-1)(p-5)/12$ (obtained by differentiating
$\Phi_p(x)=1+\cdots+x^{p-1}$ at $x=1$) give
$\Tr f_{J,1}=-J(p^2-1)/6$, hence \eqref{eq:blind-clifford} at $q=1$; the
case $q=-1$ is the sign \eqref{eq:lens-symmetries}.  For the rank
increment, $C_{J+p}=C_J$ at $p$th roots of unity while
$C'_{J+p}-C'_J=p(x^{J+1}-x^{-J})/(x-1)$, and substitution gives
$f_{J+p,1}-f_{J,1}=2px/(x-1)^2$, whose trace is $-p(p^2-1)/6$.
\end{proof}

\begin{remark}[Cyclotomic-trace route]
\label{rem:trace-route}
Theorem~\ref{thm:E} was first proved by a different route, which we
sketch because it explains the shape of the congruence.  By
Lemma~\ref{lem:trace-form}, for $q\not\equiv\pm1$ one has
$p\,\eta_J=\Tr_{K/\Q}f(\zeta)$ with $K=\Q(\zeta)$, $\zeta=\e^{2\pi i/p}$,
$f=C\,[A_q+A_{\bar q}]$ as in \eqref{eq:C-A} ($\bar q$ a positive integer
inverse of $q$), and $p\,\eta_J=\pm\Tr f_{J,1}(\zeta)$ for $q\equiv\pm1$.
The denominators factor as
$x^r+x^{-r}-x-x^{-1}=x^{-r}(x^{r-1}-1)(x^{r+1}-1)$, so $f$ is integral
away from $p$, and at $\lambda=1-\zeta$ it has a pole of order at most
two with $p$-integral Laurent coefficients:
$f(\zeta)=c_{-2}\lambda^{-2}+c_{-1}\lambda^{-1}+c_0+\lambda h$,
$h\in\Z_{(p)}[\zeta]$.  Since $\Tr(\lambda h)\in p\Z_{(p)}$ and the
traces of $\lambda^{-1},\lambda^{-2},1$ are congruent modulo $p$ to
$-\tfrac12,-\tfrac5{12},-1$, which are minus the constant terms of
$(1-\e^u)^{-1},(1-\e^u)^{-2},1$, one gets
$\Tr g(\zeta)\equiv-[u^0]g(\e^u)\pmod p$ for such $g$.  The constant
term of $f(\e^u)$ is
$4J\sum_{r\in\{q,\bar q\}}\frac r{r^2-1}\bigl[\frac{J^2}2+\frac J2+\frac{J^2r^2}6-\frac{r^2}{12}\bigr]$,
which after $\bar q\equiv q^{-1}$ collapses to $J(2J^2-1)(q^2+1)/(3q)$;
this is (ii).  Evenness follows because $f$ is even under $x\mapsto1/x$,
so $\Tr_{K/\Q}f=2\Tr_{K^+/\Q}f$, and $f$ lies in the inverse different of
$K^+=\Q(\zeta+\zeta^{-1})$ (valuation $\geq-1$ at the prime above $p$,
against the exponent $-(p-3)/2$ for $p\geq5$).
\end{remark}

\begin{remark}
Even integrality never by itself gives a zero: $p=7$, $J=2$, $q=1$ is a
silver row with $\eta_2(L(7,1))=-2$.  The boundary $p\geq5$ is
substantive: $\eta_1(L(3,1))=-2/9$, so (i) fails at $p=3$, and the
division by $3$ in (ii) is unavailable there.
\end{remark}

\subsection{Every modulus}

\begin{theorem}[Universal congruence]
\label{thm:U}
For every modulus $p\geq2$, every unit $q$ modulo $p$ and every $J\geq1$,
\begin{equation}
 \eta_J(L(p,q))-J(2J^2-1)\,\eta_1(L(p,q))\in2\Z,\qquad
 \eta_1(L(p,q))=-4s(q,p).
 \label{eq:U}
\end{equation}
\end{theorem}

\begin{proof}
Put $N=J(J^2-1)/6=\sum_{a=1}^{J-1}a(J-a)\in\Z$ and fix representatives
$q,q^*\in[1,p-1]$ with $qq^*\equiv1$.  Subtracting
$J(2J^2-1)\eta_1=-4J(2J^2-1)s$ from \eqref{eq:master}, the Dedekind term
becomes $-4Js+4J(2J^2-1)s=8J(J^2-1)s=48Ns$.  For the residue term use
$\rho_p(k)=p-2\langle k\rangle_p-p\,\delta_p(k)$ and
$\langle aq\rangle_p=aq-p\lfloor aq/p\rfloor$:
\[
 \frac2p\sum_{a=1}^{J-1}(J-a)\bigl[\rho_p(aq)+\rho_p(aq^*)\bigr]
 =-\frac{4N}p(q+q^*)+(\text{even integers}).
\]
The aliasing double sum in \eqref{eq:master} is $2\times$ an integer, and
$L_J$ is even by Lemma~\ref{lem:low-mode-parity}.  Hence
\[
 \eta_J-J(2J^2-1)\eta_1
 =\frac{4N}p\bigl(12ps-(q+q^*)\bigr)+(\text{even integers})-L_J .
\]
By \eqref{eq:rademacher-congruence}, $12ps-(q+q^*)=T-(q+q^*)\in p\Z$, so
the first term lies in $4N\Z\subseteq2\Z$.
\end{proof}

\begin{theorem}[Exact denominator and necessary conditions for a zero]
\label{thm:denominator}
Let $p\geq2$, $q$ a unit, $J\geq1$, and put $C=J(2J^2-1)$ and
$T=12p\,s(q,p)\in2\Z$.  Then
\begin{enumerate}[label=(\alph*)]
\item the reduced denominator of $\eta_J(L(p,q))$ is
$3p/\gcd(3p,\,CT)$; in particular $\eta_J\in\Z$ iff $3p\mid CT$, and
$\eta_J\in2\Z$ iff $6p\mid CT$;
\item if $\eta_J(L(p,q))=0$ then $p\mid C(q^2+1)$; if moreover $4\mid p$
then $2^{v_2(p)}\mid J$, and if $3\mid p$ then $3^{v_3(p)+1}\mid J$;
\item writing $p=2^a3^bm$ with $\gcd(m,6)=1$,
\[
 \eta_J(L(p,q))\in2\Z\iff
 m\mid C(q^2+1)\ \text{ and }\ (a<2\text{ or }2^a\mid J)\ \text{ and }\
 (b=0\text{ or }3^{b+1}\mid J);
\]
\item if $3\nmid p$, then $p\,\eta_J\in2\Z$ and
$p\,\eta_J\equiv-C(q+q^*)/3\pmod p$; that is, both parts of
Theorem~\ref{thm:E} hold at every modulus coprime to $3$.
\end{enumerate}
\end{theorem}

\begin{proof}
By Theorem~\ref{thm:U} and $\eta_1=-T/(3p)$ there is $H\in\Z$ with
$\eta_J=-CT/(3p)+2H$; adding an integer does not change a reduced
denominator, which proves (a).  For (b), a zero gives $3p\mid CT$, hence
$p\mid CT$, and by \eqref{eq:rademacher-congruence}
$C(q^2+1)\equiv qCT\equiv0\pmod p$.  If $4\mid p$ then $q$ is odd,
$q^2+1\equiv2\pmod4$, and \eqref{eq:rademacher-congruence} modulo $4$
gives $v_2(T)=1$; a zero needs $6p\mid CT$, so $v_2(C)\geq v_2(p)$, and
$v_2(C)=v_2(J)$ since $2J^2-1$ is odd.  If $3\mid p$ then
\eqref{eq:rademacher-congruence} modulo $3$ gives $qT\equiv q^2+1\equiv2$,
so $v_3(T)=0$, and $6p\mid CT$ forces $v_3(C)\geq v_3(p)+1$;
$v_3(C)=v_3(J)$ because $2J^2-1\not\equiv0\pmod3$.  For (c), the
$m$-part: $T\equiv q+q^*\pmod m$, so $m\mid CT$ iff $m\mid C(q+q^*)$ iff
$m\mid C(q^2+1)$.  The $2$-part: $T$ is always even; if $a=1$, write
$p=2u$ with $u$ odd: $r_k\equiv k\pmod2$ so $R\equiv\sum k=u(2u-1)$ is
odd, $u\mid3R$ makes $3R/u$ odd, and
$T=2[3R/u-3u(2u-1)]$ is divisible by $4$; if $a\geq2$, $v_2(T)=1$ as
above.  The $3$-part: if $b=0$ then $p\mid6R$ and $\gcd(p,3)=1$ give
$p\mid2R$, so $3\mid T$; if $b\geq1$, $v_3(T)=0$.  Comparing with
$v_2(6p)=1+a$ and $v_3(6p)=1+b$ gives (c).  For (d), if $3\nmid p$ then
$T\in6\Z$ by the $3$-part just shown, so $CT/3\in2\Z$ and
$p\,\eta_J=-CT/3+2pH\in2\Z$; reducing modulo $p$ and using
$T\equiv q+q^*$ with $3$ invertible gives the congruence.
\end{proof}

\begin{remark}[Boundaries]
\label{rem:U-boundaries}
The conditions in (b) are necessary, not sufficient:
$\eta_{25}(L(25,1))=-104$ although $25\mid C(q^2+1)$ with
$C=25(2\cdot625-1)$.  At even moduli an integral value need not be
even: $\eta_2(L(4,1))=-1$, $\eta_1(L(4,1))=-1/2$, and the difference
\eqref{eq:U} is $6$.  At $p=2$ the two branch multiplicities agree at
every level, so $\eta_J(L(2,1))=0$ for all $J$ and (d) holds trivially.
Prime powers must be read with valuations: at $p=25$, $J=5$, $q=2$ one
factor $5$ comes from $J$ and one from $q^2+1$.  Part (b) with the full
modulus replaced by its largest divisor coprime to $6$ is the weaker
statement one obtains by localizing at primes $\ell\geq5$.
\end{remark}

\begin{remark}[Status of the zero problem after this section]
\label{rem:zero-status}
For prime $p\geq5$ the zero set of $\eta_J(L(p,\cdot))$ is known at every
rank with $2J^2\not\equiv1\pmod p$ (Corollary~\ref{cor:E}).  At a silver
rank it is an even integer whose zeros include the balanced family (b),
the family of Theorem~\ref{thm:C3} below, mirror spaces when
$p\equiv1\pmod4$, and further zeros not yet characterized
(Section~\ref{sec:silver-zero-set}).  At composite moduli
Theorem~\ref{thm:denominator} gives necessary local conditions only;
Observation~\ref{obs:composite} shows that cross-cancellations between
primitive parts occur.  Congruences of the type of
Theorems~\ref{thm:E} and \ref{thm:U} may follow from index theory for
flat bundles; that literature has not been searched, and no novelty is
claimed for them.
\end{remark}

\section{The silver zero at \texorpdfstring{$q=J$}{q = J}}
\label{sec:silver}

\subsection{Statement}

\begin{theorem}
\label{thm:C3}
For every $J\geq2$ and every divisor $p>1$ of $2J^2-1$,
\[
 \eta_J(L(p,J))=0,
\]
and hence also $\eta_J(L(p,q))=0$ for $q\equiv\pm J$ and $q\equiv\pm2J
\pmod p$.
\end{theorem}

For $p=2J^2-1$ this is the first family of
\cite[Thm.~7.1]{MahmoodP1Collection2026}, proved there by repeated
Dedekind reciprocity with residues that do not wrap modulo $p$.  At a
proper divisor the residues wrap and that argument does not apply; the
proof below uses instead the conjugacy invariance of Rademacher's symbol.
The transfer to $q\equiv\pm J,\pm2J$ is \eqref{eq:lens-symmetries}
together with $2J\equiv J^{-1}\pmod p$.

Throughout, fix $J\geq2$ and a divisor $p>1$ of $N:=2J^2-1$.  Then $p$
is odd, coprime to $3$, and $m_0:=N/p$ is odd.  The residues $2J$,
$J\pm1$, $2J\pm1$ are units modulo $p$: a prime dividing $J\pm1$ and $N$
would divide $2J^2-1-2(J\mp1)(J\pm1)=1$; $(2J+1)(2J-1)=2N+1\equiv1$; and
$(J+1)(J-1)\equiv-\tfrac12$.  In particular $q=J$ is nonsingular in the
sense of Lemma~\ref{lem:trace-form}: no nonidentity element is a
Clifford translation, and this holds at composite $p$ as well.  Put
\begin{equation}
 \varepsilon:=4J+3,\qquad\lambda:=-\varepsilon,\qquad
 \lambda^2+6\lambda+1=\varepsilon^2-6\varepsilon+1=8N\equiv0\pmod p;
 \label{eq:silver-unit}
\end{equation}
$\varepsilon$ is the silver unit $3+2\sqrt2$ read modulo $p$ with
$\sqrt2\equiv2J$.  Two notational conventions are essential: $c_m(z)$
denotes the centered residue of $z$ modulo $m$ in $(-m/2,m/2]$; and a
fraction $a/b$ written inside $c_p(\cdot)$ means the modular quotient
$ab^{-1}$ in $\Z/p$, not a real number.  For instance
$c_p(j+\tfrac12)=j-\tfrac{p-1}2$ for $0\leq j<p$, and
$c_{17}(\tfrac12)=c_{17}(9)=-8$.  Finally $s_n(a,b;p)$ is the
Fourier--Dedekind sum \cite{BeckRobins2004}
\[
 s_n(a,b;p)=\frac1p\sum_{\substack{\zeta^p=1\\\zeta\neq1}}
 \frac{\zeta^n}{(\zeta^a-1)(\zeta^b-1)},\qquad\gcd(a,p)=\gcd(b,p)=1,
\]
which satisfies $s_n=s_{a+b-n}$ and $s_n(a,-b)=-s_{n+b}(a,b)$.

\subsection{Reduction to a quarter-shifted covariance}

\begin{lemma}[Step 1: a trigonometric sum]
\label{lem:C3-step1}
With $\varphi_k=\pi k/p$,
\[
 p\,\eta_J(L(p,J))=\sum_{k=1}^{p-1}
 \frac{(-1)^k\cos(\pi k/p)}{\cos(2\pi k/p)-\cos(2\pi Jk/p)} .
\]
\end{lemma}

\begin{proof}
Apply \eqref{eq:E1} with $q=J$.  Write $2J^2=1+m_0p$ with $m_0$ odd;
then $\sin(2J^2\varphi_k)=(-1)^k\sin\varphi_k$,
$\cos((2J+1)J\varphi_k)=(-1)^k\cos((J+1)\varphi_k)$,
$\sin(2J\varphi_k)/\sin(J\varphi_k)=2\cos(J\varphi_k)$, and
$\cos((2J+1)\varphi)-2\cos((J+1)\varphi)\cos(J\varphi)=-\cos\varphi$.
The numerator of \eqref{eq:E1} becomes $-(-1)^k\cos\varphi_k$ and the
denominator $\cos(2J\varphi_k)-\cos(2\varphi_k)$.
\end{proof}

\begin{lemma}[Step 2: Fourier--Dedekind form]
\label{lem:C3-step2}
$p\,\eta_J(L(p,J))=2p\,s_J(2J+1,\,2J-1;\,p)$.
\end{lemma}

\begin{proof}
Put $h=(p+1)/2$, $\theta=2\pi k/p$ and $\zeta=\e^{i\theta}$.  Then
$(-1)^k\cos(\pi k/p)=\cos(h\theta)$ and
$\cos\theta-\cos(J\theta)=-\tfrac12\zeta^{-J}(\zeta^{J+1}-1)(\zeta^{J-1}-1)$,
so by Lemma~\ref{lem:C3-step1}, using $s_n=s_{a+b-n}$ with $a+b=2J$,
$p\,\eta_J=-2p\,s_{J+h}(J+1,J-1;p)$.  The substitution
$\zeta\mapsto\zeta^{2J}$ (a permutation of the nontrivial $p$th roots of
unity, as $\gcd(2J,p)=1$) sends the exponents $(J+1,J-1,J+h)$ to
$(2J+1,-(2J-1),J+1)$ modulo $p$, because $2J^2\equiv1$ and $2h\equiv1$.
Then $s_n(a,-b)=-s_{n+b}(a,b)$ and once more $s_n=s_{a+b-n}$ give
$s_{J+h}(J+1,J-1)=-s_{3J}(2J+1,2J-1)=-s_J(2J+1,2J-1)$.
\end{proof}

\begin{lemma}[Step 3: covariance form]
\label{lem:C3-step3}
Let $K:=\sum_{x\in\Z/p}c_p\bigl(x+\tfrac14\bigr)\,c_p\bigl(\lambda x+\tfrac14\bigr)$.
Then $\eta_J(L(p,J))=2K/p^2$; in particular
$\eta_J(L(p,J))=0\iff K=0$.
\end{lemma}

\begin{proof}
For $\zeta^p=1$, $\zeta\neq1$ and $a$ a unit,
$1/(\zeta^a-1)=\frac1p\sum_{j=0}^{p-1}j\zeta^{aj}$.  Inserting this for
both factors of $s_n(a,b;p)$ and using orthogonality of the $p$th roots
of unity gives, with $\mu=(p-1)/2$,
\[
 p\,s_n(a,b;p)=\frac1p\sum_{j_1=0}^{p-1}(j_1-\mu)(j_2-\mu),
 \qquad aj_1+bj_2+n\equiv0\pmod p,\ 0\leq j_2<p .
\]
Now $j-\mu=c_p(j+\tfrac12)$.  For $(a,b,n)=(2J+1,2J-1,J)$, using
$b^{-1}\equiv a$ and $a^2\equiv4J+3=\varepsilon$, the constraint reads
$y_2\equiv-\varepsilon y_1+(1+J)$ for $y_i=j_i+\tfrac12$, and
$1+J=(1+\varepsilon)/4$ in $\Z/p$.  With $x=y_1-\tfrac14$ this is
$y_2=\lambda x+\tfrac14$, so $p\,s_J=K/p$, and Lemma~\ref{lem:C3-step2}
gives $\eta_J=2K/p^2$.
\end{proof}

\subsection{From the quarter shift to Dedekind sums at even moduli}

\begin{lemma}[Step 4: lifting to modulus $4p$]
\label{lem:C3-step4}
Put $u=c_p(x+\tfrac14)$, $w=c_p(\lambda x+\tfrac14)$, $U=4u-1$,
$W=4w-1$.  Fix $\hat\lambda\in\Z$ with $\hat\lambda\equiv\lambda\pmod p$
and $\hat\lambda\equiv1\pmod4$.  Then $\hat\lambda^2+6\hat\lambda+1\equiv0
\pmod{4p}$, and as $x$ runs over $\Z/p$, $U$ runs over the $p$ integers
$\equiv3\pmod4$ in $(-2p,2p)$, with $W=c_{4p}(\hat\lambda U)$.
Consequently
\begin{equation}
 16K=T-p,\qquad
 T:=\sum_{\substack{U\in\Z/4p\\U\equiv3\,(4)}}c_{4p}(U)\,c_{4p}(\hat\lambda U).
 \label{eq:16K}
\end{equation}
\end{lemma}

\begin{proof}
The congruence modulo $4p$ holds modulo $p$ by \eqref{eq:silver-unit}
and modulo $4$ because $1+6+1=8$.  As $x$ runs over $\Z/p$ so does
$x+\tfrac14$, so $u$ runs over the centered residues and $U=4u-1$ over the
stated set; since $\lambda$ is a unit the same holds for $w$ and $W$.
Now $W\equiv4\lambda x\equiv\hat\lambda U\pmod p$ and
$W\equiv3\equiv\hat\lambda U\pmod4$, so $W\equiv\hat\lambda U\pmod{4p}$,
and $W\in(-2p,2p)$ is the centered representative.  Finally
$\sum U=\sum W=-p$ (the centered residues sum to zero), so
$16K=\sum(U+1)(W+1)=T-p-p+p$.
\end{proof}

\begin{lemma}[Step 5: parity split]
\label{lem:C3-step5}
For an even modulus $m$ and $h$ odd and coprime to $m$ put
$D_m(h):=\sum_{V\in\Z/m}c_m(V)\,c_m(hV)$.  Then
$2T=D_{4p}(\hat\lambda)-4D_{2p}(\hat\lambda)$.
\end{lemma}

\begin{proof}
$U\mapsto-U$ exchanges the odd classes $1$ and $3$ modulo $4$ and
preserves each summand ($\hat\lambda$ is odd and $\pm2p$ is even), so the
sum over odd $U$ is $2T$.  The even $U=2V$ contribute
$c_{4p}(2V)c_{4p}(2\hat\lambda V)=4c_{2p}(V)c_{2p}(\hat\lambda V)$.
\end{proof}

\begin{lemma}[Centered residues and Dedekind sums]
\label{lem:D}
For even $m$ and $h$ odd and coprime to $m$,
$D_m(h)=\dfrac{m^2}4+m^2\,s(h,m)$.
\end{lemma}

\begin{proof}
For $V\not\equiv m/2$, $c_m(V)/m=\sawtooth{V/m+\tfrac12}$; at $V\equiv m/2$
the left side is $\tfrac12$ and the right side $0$, and $hV\equiv m/2$
iff $V\equiv m/2$ because $h$ is odd.  Hence
$D_m/m^2=\tfrac14+\sum_V\sawtooth{V/m+\tfrac12}\sawtooth{hV/m+\tfrac12}$.
Use $\sawtooth{x+\tfrac12}=\sawtooth{2x}-\sawtooth x$ in both factors.
Of the four resulting sums, $\sum_V\sawtooth{2V/m}\sawtooth{2hV/m}=2s(h,m/2)$;
the two mixed sums each equal $s(h,m/2)$, by splitting $V$ modulo $m/2$
and using $\sawtooth y+\sawtooth{y+\tfrac12}=\sawtooth{2y}$; and
$\sum_V\sawtooth{V/m}\sawtooth{hV/m}=s(h,m)$.  The signed total is
$2s(h,m/2)-2s(h,m/2)+s(h,m)$.
\end{proof}

Combining Lemmas~\ref{lem:C3-step4}--\ref{lem:D},
$2T=16p^2\bigl[s(\hat\lambda,4p)-s(\hat\lambda,2p)\bigr]$, and therefore
\begin{equation}
 \eta_J(L(p,J))=0\iff K=0\iff T=p\iff
 s(\hat\lambda,4p)-s(\hat\lambda,2p)=\frac1{8p}.
 \label{eq:C3-reduced}
\end{equation}

\subsection{Rademacher's symbol and the centralizer}

For $M=\begin{pmatrix}a&b\\c&d\end{pmatrix}\in SL_2(\Z)$ Rademacher's
functions are
\begin{equation}
 \Phi(M)=\begin{cases}\dfrac{a+d}c-12\,\sgn(c)\,s(d,|c|),&c\neq0,\\[6pt]
 \dfrac bd,&c=0,\end{cases}
 \qquad
 \Psi(M)=\Phi(M)-3\,\sgn\bigl(c(a+d)\bigr),
 \label{eq:rademacher-symbol}
\end{equation}
and $\Psi$ is invariant under $SL_2(\Z)$-conjugation on hyperbolic
elements \cite[Ch.~4]{RademacherGrosswald1972}; this normalization
agrees with \cite[eq.~(1.2)]{Burrin2022} and
\cite[eq.~(1.2)]{Matsusaka2024}.  Directly from the definition, with
$F=\operatorname{diag}(1,-1)$,
\begin{equation}
 \Psi(FMF)=-\Psi(M),
 \label{eq:psi-reflection}
\end{equation}
since $FMF=\begin{pmatrix}a&-b\\-c&d\end{pmatrix}$ negates $\Phi$ and
$\sgn(c(a+d))$.

\begin{lemma}[Dedekind sums at the silver unit, even moduli]
\label{lem:R}
For $m\in\{2p,4p\}$, $s(\hat\lambda,m)=\dfrac{m-2}{4m}$.
\end{lemma}

\begin{proof}
Write $\hat\lambda=1+4t$; then $\hat\lambda^2+6\hat\lambda+1=8\sigma$
with $\sigma=1+4t+2t^2$ odd and $p\mid\sigma$.  For $m\in\{2p,4p\}$ the
integer $b:=-(\hat\lambda^2+6\hat\lambda+1)/m$ is even ($-4\sigma/p$ or
$-2\sigma/p$).  Let
\[
 M=\begin{pmatrix}-6-\hat\lambda&b\\ m&\hat\lambda\end{pmatrix}\in SL_2(\Z),
 \qquad\operatorname{tr}M=-6,
\]
which is hyperbolic.  Since $\hat\lambda+3\equiv0\pmod4$ and $b,m$ are
even, $A:=(M+3I)/2$ is integral; Cayley--Hamilton ($M^2+6M+I=0$) gives
$A^2=2I$, $\operatorname{tr}A=0$, $\det A=-2$.  Hence $R:=I+A$ is
integral, commutes with $M$, and has $\det R=(1+\sqrt2)(1-\sqrt2)=-1$.
Put $P=RF$; then $\det P=1$ and $P^{-1}MP=FR^{-1}MRF=FMF$.  By
conjugation invariance and \eqref{eq:psi-reflection},
$\Psi(M)=\Psi(FMF)=-\Psi(M)$, so $\Psi(M)=0$.  Now $c=m>0$ and
$a+d=-6<0$, so
\[
 0=\Psi(M)=-\frac6m-12\,s(\hat\lambda,m)+3,\qquad
 s(\hat\lambda,m)=\frac{3-6/m}{12}=\frac{m-2}{4m}.
\]
With the unreduced lift $\hat\lambda=-(4J+3)$ (already $\equiv1\pmod4$)
and $n_0=N/p$, the centralizers are explicit:
$A_{2p}=\begin{pmatrix}2J&-2n_0\\p&-2J\end{pmatrix}$ and
$A_{4p}=\begin{pmatrix}2J&-n_0\\2p&-2J\end{pmatrix}$, both squaring to
$2I$ because $4J^2-2pn_0=2$.
\end{proof}

\begin{proof}[Proof of Theorem~\ref{thm:C3}]
By Lemma~\ref{lem:R},
$s(\hat\lambda,4p)-s(\hat\lambda,2p)=\frac{4p-2}{16p}-\frac{2p-2}{8p}=\frac1{8p}$,
so \eqref{eq:C3-reduced} gives $K=0$ and $\eta_J(L(p,J))=0$.
\end{proof}

\begin{remark}[The binary-form route]
\label{rem:class-number-route}
Lemma~\ref{lem:R} was first proved through binary quadratic forms.  The
form $Q_M=[m,\,2\hat\lambda+6,\,-b]$ attached to $M$ (via
$M\mapsto[c,\,d-a,\,-b]$) has discriminant $32$ and content exactly $2$,
since $32/\content^2$ must be $\equiv0,1\pmod4$.  Its primitive part has
discriminant $8$; because $\Z[\sqrt2]$ is norm-Euclidean, hence a
principal ideal domain, and has the unit $1+\sqrt2$ of norm $-1$, there
is a single proper class of primitive forms of discriminant $8$, so $M$
is $SL_2(\Z)$-conjugate to $M_0=\begin{pmatrix}-3&4\\2&-3\end{pmatrix}$
(whose form is $[2,0,-4]$), and $\Psi(M_0)=-3-12s(1,2)+3=0$.  The
$SL_2$ (rather than $GL_2$) conjugacy is obtained by adjusting an
integral $O$-basis of $\Z^2$ by $I+A_0$, $A_0=\begin{pmatrix}0&2\\1&0\end{pmatrix}$,
which has determinant $-1$ and commutes with $M_0$.  The centralizer
proof above avoids the class-number input.  Both routes rest on the
conjugacy invariance of $\Psi$, which we cite and do not re-prove.
\end{remark}

\begin{remark}[Why a quarter shift, and why silver ranks]
The shift $\tfrac14$ in Lemma~\ref{lem:C3-step3} is what forces
$\hat\lambda\equiv1\pmod4$.  With $\hat\lambda\equiv3\pmod4$ the form at
$m=4p$ is primitive, $\Psi=\pm3$, and the evaluation fails (for
$(J,p)=(3,17)$ one gets $s=-1/136$ instead of $33/136$).  The unit
$1+\sqrt2$ of norm $-1$ is what makes the content-two class have
$\Psi=0$; in the language of Section~\ref{sec:rademacher}, $M$ is
\emph{inert}.
\end{remark}

\subsection{Consequences and clarifications}

\begin{corollary}[The Dedekind sum at the silver unit, odd moduli]
\label{cor:S}
Let $p$ be odd and let $\varepsilon\in\Z/p$ satisfy
$\varepsilon^2-6\varepsilon+1\equiv0\pmod p$.  Then
\[
 s(\varepsilon,p)=
 \begin{cases}-\dfrac{p-1}{2p},&p\equiv1\pmod8,\\[6pt]
 \dfrac1{2p},&p\equiv7\pmod8.\end{cases}
\]
\end{corollary}

\begin{proof}
Lift $h=-\varepsilon$ to $h\equiv1\pmod4$ and write
$h^2+6h+1=8\sigma$ with $\sigma$ odd and $p\mid\sigma$.  The matrix
$M_p=\begin{pmatrix}-6-h&-8\sigma/p\\p&h\end{pmatrix}$ has trace $-6$,
and its form $H=[p,\,2h+6,\,8\sigma/p]$ has discriminant $32$ and is
primitive ($p$ is odd).  The integral form $G=[p,\,h+3,\,2\sigma/p]$ has
discriminant $8$ and $H(x,y)=G(x,2y)$; by Remark~\ref{rem:class-number-route}
$G$ is properly equivalent to $x^2-2y^2$, so $H$ is that norm form
restricted to an oriented sublattice of index two.  There are three such
sublattices, the preimages of the three lines in $(\Z/2)^2$: on
$\{x\text{ even}\}$ the norm is always even, impossible since $H$
represents the odd number $p$; on $\{y\text{ even}\}$ the restricted form is
$x^2-8y^2$; on $\{x\equiv y\pmod2\}$, with the positively oriented basis
$(1,1),(0,2)$, it is $-x^2-8xy-8y^2$, properly equivalent to $-x^2+8y^2$
by a unimodular substitution.  The odd values of these two forms are
respectively $\equiv1$ and $\equiv7\pmod8$, so they are not properly
equivalent, and $H$ is determined by $p\bmod8$.  At fixed trace $t$ the
map $M\mapsto Q_M=[c,d-a,-b]$ is injective: from $[u,v,w]$ one recovers
$c=u$, $b=-w$, $a=(t-v)/2$ and $d=(t+v)/2$.  It is equivariant because
$Q_M(x)=\det(x,Mx)$ gives $Q_{P^{-1}MP}(x)=Q_M(Px)$ for
$P\in SL_2(\Z)$.  Thus $M_p$ is conjugate to
$M_+=\begin{pmatrix}-3&8\\1&-3\end{pmatrix}$ ($\Psi=-3$) if
$p\equiv1\pmod8$ and to $M_-=\begin{pmatrix}-3&-8\\-1&-3\end{pmatrix}$
($\Psi=+3$) if $p\equiv7\pmod8$.  Since
$\Psi(M_p)=3-6/p-12s(h,p)=3-6/p+12s(\varepsilon,p)$, solving gives the
two values.
\end{proof}

\begin{proposition}[The silver zero and level balance]
\label{prop:C3-balance}
Let $J\geq2$ and $p>1$ divide $2J^2-1$.  Then $L(p,J)$ is chiral, and it
is rank-$J$ balanced if and only if $p=7$.  Thus for $p\neq7$
Theorem~\ref{thm:C3} produces zeros of $\eta_J$ which are not level-wise
balanced; for example at $(J,p)=(3,17)$ the imbalance sequence begins
$D_3=(0,0,-2,4,-2,2,-4,2,0,\ldots)$, and only its Hurwitz moments cancel.
\end{proposition}

\begin{proof}
Mirror balance would require $J^2\equiv-1$ together with $2J^2\equiv1$,
forcing $p\mid3$, but $3\nmid2J^2-1$.  A blind rank is impossible since
$J$ is a unit.  For the silver family, $J\equiv2J\pm1$ forces $J\equiv\mp1$,
contradicting $p\mid2J^2-1$; $J\equiv-(2J\pm1)$ gives $3J\equiv\mp1$,
hence $9J^2\equiv1$ and, with $2J^2\equiv1$, $9\equiv2$, i.e.\ $p\mid7$.
Conversely at $p=7$, $2J^2\equiv1$ gives $J\equiv\pm2$, and
$q=J$ then lies in $\{\pm(2J\pm1)\}=\{2,3,4,5\}$.  The sequence quoted is
an exact computation.
\end{proof}

\begin{remark}
The other universal zero orbit at silver ranks, $q\equiv\pm(2J\pm1)$, is
the balanced family (b) of Theorem~\ref{thm:classification}; in
\eqref{eq:E1} its two terms are exchanged with opposite signs by
$k\mapsto(2J+1)k$, which gives a second, Fourier-side proof of
$\eta_J=0$ for that family.  Theorems~\ref{thm:E} and \ref{thm:U} are
not used in the proof of Theorem~\ref{thm:C3}; they are consistent with
it, an exact vanishing being finer than any congruence.
\end{remark}

\section{Zero versus symmetry: the Rademacher symbol}
\label{sec:rademacher}

The proof of Lemma~\ref{lem:R} turned on an integral matrix of
determinant $-1$ commuting with a hyperbolic $M\in SL_2(\Z)$, which forces
$\Psi(M)=0$.  This section records, with citations, the extent to which
that mechanism is necessary.  None of the facts about $\Psi$ below is
claimed as new: the forward implications are classical, and the negative
answer also follows from known asymptotic counts.  What we add are
explicit small witnesses with elementary proofs, exact statements for the
two- and four-block classes, an elementary density bound, and the
identification of the matrices behind Theorem~\ref{thm:C3}.

\subsection{Conventions and cited inputs}

Let $M=\begin{pmatrix}a&b\\c&d\end{pmatrix}\in SL_2(\Z)$ be hyperbolic
with $t=a+d$, and $\Phi,\Psi$ as in \eqref{eq:rademacher-symbol}.  Put
$R=\begin{pmatrix}1&1\\0&1\end{pmatrix}$,
$L=\begin{pmatrix}1&0\\1&1\end{pmatrix}$, $F=\operatorname{diag}(1,-1)$,
$S=\begin{pmatrix}0&-1\\1&0\end{pmatrix}$ and
$W=\begin{pmatrix}0&1\\1&0\end{pmatrix}=SF$.  The form
$Q_M=[c,\,d-a,\,-b]$ has content $g$, and $Q_M/g=[A',B',C']$ has
discriminant $D=(t^2-4)/g^2$; $O_D$ is the order of discriminant $D$.
For $t>2$, $w(M)$ denotes the cyclic word in $R,L$ of $M$ (input (K3)
below); $\sigma$ swaps the two letters, $\rho$ reverses the word, and
$w\sim w'$ means equality up to rotation.  A class is \emph{primitive}
if $M$ is not a proper power, i.e.\ its word is aperiodic.  Following
Sarnak \cite{Sarnak2007}, $M$ is \emph{inert} if $M\sim FMF$,
\emph{reciprocal} if $M\sim M^{-1}$, and \emph{ambiguous} if
$M\sim FM^{-1}F$, all conjugacies in $SL_2(\Z)$.  A zero $\Psi(M)=0$ is
called \emph{accidental} if $M$ is neither inert nor reciprocal.  Classes
with $t<-2$ are handled through $-M$: $\Psi(-M)=\Psi(M)$, and the three
symmetries commute with negation.

The following inputs are cited, not re-proved.
\begin{enumerate}[label=(K\arabic*)]
\item $\Psi$ is an $SL_2(\Z)$-class function on hyperbolic elements
\cite[Ch.~4]{RademacherGrosswald1972}.
\item $\Psi(M^{-1})=-\Psi(M)$ \cite[Lem.~6.1(ii)]{BurrinVonEssen2024}.
\item Every $M$ with $t>2$ is $SL_2(\Z)$-conjugate to a positive word in
$R$ and $L$ containing both letters, unique up to rotation; $M$ is
primitive iff the word is aperiodic.  This is the cyclically reduced
normal form in $PSL_2(\Z)\cong C_2*C_3$; see
\cite[pp.~268,272]{Ghys2006} for the free-product and left-right coding.
\item $\Psi(M)=\#R-\#L$, the alternating sum of the run lengths of $w(M)$
\cite[Prop.~4.1]{BurrinVonEssen2024}; see also Zagier
\cite{Zagier1975}.
\item \emph{(Context only.)}  $-\Psi(M)/3$ is the Atiyah--Patodi--Singer
eta invariant of the mapping torus $T_M$ and its adiabatic limit, in the
normalization of \cite[eq.~(5.60)]{Atiyah1987} as restated in
\cite[Sec.~4.2]{BurrinVonEssen2024}; $\Psi(M)$ is also the linking number
of the modular knot of $M$ with the trefoil \cite{Ghys2006}.  The
statement ``$T_M$ has vanishing eta invariant'' below refers to this
cited normalization only, not to arbitrary metrics or other operators.
\item For Anosov $M$, $T_M$ admits an orientation-reversing
self-homeomorphism iff there is an integral $X$ with $\det X=-1$ and
$XM=MX$, or with $\det X=+1$ and $XM=M^{-1}X$
\cite[Sec.~1.3 and the Anosov case of the degree theorem]{SunWangWu2008}.
\item Sarnak's counts by trace \cite[Thm.~2]{Sarnak2007} and the
Sarnak--Mozzochi winding law as stated in
\cite[Cor.~1.3]{BurrinVonEssen2024}; used only in
Theorem~\ref{thm:Z}(vi)(b).
\end{enumerate}
Two elementary facts: $\Psi(FMF)=-\Psi(M)$ (eq.~\eqref{eq:psi-reflection});
and $WRW=L$, $WLW=R$, $SR^{-1}S^{-1}=L$, $SL^{-1}S^{-1}=R$.

\begin{lemma}[Witness dictionary]
\label{lem:W}
Let $t>2$ and $w=w(M)$.
\begin{enumerate}[label=(\alph*)]
\item The following are equivalent: (1) there is $X\in M_2(\Z)$ with
$\det X=-1$ and $XM=MX$; (2) $P^{-1}MP=FMF$ for some $P\in SL_2(\Z)$;
(3) $\sigma(w)\sim w$; (4) $x^2-Dy^2=-4$ has an integer solution, i.e.\
$O_D$ has a unit of norm $-1$.
\item $M$ is reciprocal iff $\sigma\rho(w)\sim w$.
\item $M$ is ambiguous iff $\rho(w)\sim w$.
\end{enumerate}
\end{lemma}

\begin{proof}
(1)$\Leftrightarrow$(2): given $X$ put $P=XF$; given $P$ put $X=PF$; in
both directions $X^{-1}MX=M\iff(XF)^{-1}M(XF)=FMF$.
(2)$\Leftrightarrow$(3): $WMW=S(FMF)S^{-1}$, and $WMW$ has the word
$\sigma(w)$; apply (K3).
(1)$\Leftrightarrow$(4): write $M=\begin{pmatrix}(t-gB')/2&-gC'\\gA'&(t+gB')/2\end{pmatrix}$.
The rational centralizer of $M$ is $\Q[M]$, and its integral elements are
$X(x,y)=\begin{pmatrix}(x-B'y)/2&-C'y\\A'y&(x+B'y)/2\end{pmatrix}$ with
$x,y\in\Z$, $x\equiv B'y\pmod2$ ($y$ is integral because $[A',B',C']$ is
primitive), and $\det X(x,y)=(x^2-Dy^2)/4$.
(b) Inversion reverses the letters, and conjugation by $S$ then gives the
word $\sigma\rho(w)$ for $M^{-1}$.
(c) $FM^{-1}F\sim WM^{-1}W$, and conjugating by $S$ gives the word $\rho(w)$.
\end{proof}

\subsection{What a zero of \texorpdfstring{$\Psi$}{Psi} detects and forgets}

\begin{theorem}
\label{thm:Z}
\begin{enumerate}[label=(\roman*)]
\item \textup{(Witnessed zeros.)}  If $M$ is inert or reciprocal then
$\Psi(M)=0$.  Ambiguity alone forces nothing: $R^2L$ is ambiguous and
$\Psi(R^2L)=1$.
\item \textup{(The inert witness is not necessary.)}
$M_1=RLR^2L^2=\begin{pmatrix}12&5\\7&3\end{pmatrix}$ has $\Psi(M_1)=0$
and $D=221$; the order $O_{221}$ has no unit of norm $-1$, so no
integral matrix of determinant $-1$ commutes with $M_1$.  Its zero is
witnessed by reciprocity: $P=RLS=\begin{pmatrix}1&-2\\1&-1\end{pmatrix}$
satisfies $P^{-1}M_1P=M_1^{-1}$.
\item \textup{(Neither inert nor reciprocal witness is necessary.)}
$M_2=L^3R^2LR^2=\begin{pmatrix}3&8\\10&27\end{pmatrix}$ has
$\Psi(M_2)=0$ and $D=224$; it is neither inert nor reciprocal.  (It is
ambiguous: reversing $LLLRRLRR$ gives a rotation of it.  Ambiguity does
not negate $\Psi$ and does not force the zero.)  By (K6) the Sol
manifold $T_{M_2}$ is chiral, that is, it admits no orientation-reversing
self-homeomorphism, while by (K5) its eta invariant, in the cited
normalization, vanishes.
\item \textup{(Minimality.)}  Every primitive balanced class of word length
$\leq6$ is inert or reciprocal.  Word length $8$ has exactly $8$
primitive balanced classes: $2$ inert, $4$ reciprocal and not inert, and
$2$ accidental, the $\sigma$-pair $\{L^3R^2LR^2,\ L^2RL^2R^3\}$.
\item \textup{(Delimited converse.)}
\begin{enumerate}[label=(\alph*)]
\item For two-block classes $R^aL^b$ ($a,b\geq1$): $\Psi=a-b$, and
$\Psi=0\iff a=b\iff$ inert.  Here the converse holds exactly.
\item For primitive four-block classes $R^aL^bR^cL^d$
($a,b,c,d\geq1$, $(a,b)\neq(c,d)$): $\Psi=0\iff a+c=b+d$; no such class
is inert; the class is reciprocal iff $(a=d$ and $b=c)$ or $(a=b$ and
$c=d)$; every other solution of $a+c=b+d$ is an accidental zero.  In
particular $R^2LR^{n-1}L^n$ is an accidental zero for every $n\geq3$,
with $n=3$ the class of $M_2$.  For unrestricted four-block tuples
$(a,b,c,d)$, the class is inert iff $a=b=c=d$.
\end{enumerate}
\item \textup{(Density.)}
\begin{enumerate}[label=(\alph*)]
\item Let $B_{2n}$ be the number of primitive balanced classes of word
length $2n$ and $W_{2n}$ the number of those that are inert or
reciprocal.  Then
\[
 B_{2n}=\frac1{2n}\sum_{d\mid n}\mu(d)\binom{2n/d}{n/d}
 \sim\frac{4^n}{2n\sqrt{\pi n}},\qquad W_{2n}<2^{n+1},
\]
so $W_{2n}/B_{2n}\to0$.
\item \textup{(Cited, not computed.)}  By (K7), inert or reciprocal
primitive classes of trace $\leq x$ number $O(x)$, while primitive
classes with $\Psi=0$ number $\sim x^2/(12(\log x)^2)$; so in the
ordering by trace as well, almost every zero is accidental.
\end{enumerate}
\item \textup{(What is detected.)}  $\Psi(M)=0$ detects exactly the
balance $\#R=\#L$ of the cutting sequence (K4), equivalently zero
winding about the cusp, zero linking with the trefoil, and vanishing
eta invariant of $T_M$ in the cited normalization (K5).  The two
$\Psi$-negating symmetries, reflection of the fibre (inert) and time
reversal of the geodesic (reciprocal), force it.  It does not reveal
whether either of them is present.
\item \textup{(Transfer of the silver zero.)}  For every $J\geq2$, every
divisor $p>1$ of $2J^2-1$ and $m\in\{2p,4p\}$, the matrix $M$ of
Lemma~\ref{lem:R} satisfies $-M\sim R^2L^2$, the inert two-block class of
(v)(a), whose determinant $-1$ commuter is $\begin{pmatrix}2&1\\1&0\end{pmatrix}$,
i.e.\ $1+\sqrt2$.  The lens space $L(p,J)$ is chiral and, for $p\neq7$,
unbalanced (Proposition~\ref{prop:C3-balance}); so for $p\neq7$ its zero
has neither of the two named lens-space mechanisms (orientation
reversal, level balance), yet it is the image of a zero that is
witnessed on the Sol side.  At $p=7$ the zero is also level-wise
balanced.
\end{enumerate}
\end{theorem}

\begin{proof}
(i) is (K1), (K2) and \eqref{eq:psi-reflection}; $\Psi(R^2L)=2-1=1$ by
(K4), and $R^2L$ is ambiguous because $\rho(RRL)=LRR\sim RRL$.

(ii) $\Psi(M_1)=3-3=0$ by (K4); directly, $Q_{M_1}=[7,-9,-5]$,
$s(3,7)=-1/14$, $\Phi=15/7+6/7=3$ and $\Psi=0$.  A unit $\eta>1$ of
$O_{221}$ is $(x+y\sqrt{221})/2$ with $x,y\geq1$.  If $y\geq2$ then
$\eta>15>\epsilon:=(15+\sqrt{221})/2$; if $y=1$ then $x^2=221\pm4$
forces $x=15$ and norm $+1$.  So $\epsilon$ is the least unit $>1$,
every unit is $\pm\epsilon^k$, and all have norm $+1$; by
Lemma~\ref{lem:W}(a) $M_1$ is not inert.  Direct multiplication gives
$P^{-1}M_1P=\begin{pmatrix}3&-5\\-7&12\end{pmatrix}=M_1^{-1}$.

(iii) $\Psi(M_2)=4-4=0$; $Q_{M_2}=[10,24,-8]$, content $2$, $D=224$.
Not inert: $x^2-224y^2=-4$ would give $x^2\equiv3\pmod7$, a
non-residue.  Not reciprocal: $\sigma\rho(LLLRRLRR)=LLRLLRRR$ has cyclic
$R$-runs of lengths $\{1,3\}$, whereas $LLLRRLRR$ has $R$-runs $\{2,2\}$,
so the former is not a rotation of the latter.
Ambiguous: $\rho(LLLRRLRR)=RRLRRLLL$ is the rotation by three places.
The last two sentences are (K6) and (K5).

(iv) A primitive balanced word of length $\leq6$ is a rotation of
$RL$, $R^2L^2$, $RLRL$ (not primitive), $R^3L^3$, $R^2LRL^2$, $R^2L^2RL$,
or $RLR^2L^2$ and their letter swaps; each is inert or reciprocal by
Lemma~\ref{lem:W}.  At length $8$, classifying by cyclic run count
gives one two-run class ($R^4L^4$, inert), four primitive four-run
classes (two reciprocal, two accidental) and three six-run classes (one
inert, two reciprocal); the accidental pair is $\{L^3R^2LR^2,L^2RL^2R^3\}$.
Proper powers inherit the witness of their primitive root, so the first
accidental length is $8$ even among all balanced words.

(v)(a) $R^aL^b=\begin{pmatrix}1+ab&a\\b&1\end{pmatrix}$, $\Psi=a-b$ by
(K4), and $\sigma(R^aL^b)=R^bL^a\sim R^aL^b$ iff $a=b$.
(b) Primitivity of a four-block word is exactly $(a,b)\neq(c,d)$.  The
$R$-first exponent sequences of one necklace are the rotations by an
even number of places, i.e.\ $(a,b,c,d)$ and $(c,d,a,b)$.  Now
$\sigma(w)\sim R^bL^cR^dL^a$ has sequence $(b,c,d,a)$, which equals one of
the two only if $a=b=c=d$, excluded by primitivity (and giving the
unrestricted statement); $\sigma\rho(w)=R^dL^cR^bL^a$ has sequence
$(d,c,b,a)$, which equals $(a,b,c,d)$ iff $a=d$, $b=c$, and equals
$(c,d,a,b)$ iff $a=b$, $c=d$; both patterns satisfy $a+c=b+d$.  For
$(2,1,n-1,n)$: $a+c=b+d=n+1$, $a=d$ would force $n=2$, $a\neq b$, and
$(2,1)\neq(n-1,n)$.

(vi)(a) The count of aperiodic necklaces with $n$ letters of each kind
is the displayed M\"obius sum.  If $\sigma(w)=\mathrm{rot}_k(w)$ with
$k\not\equiv0$, then $w_{i+k}=1-w_i$, so $w$ is determined by its values
on the $\gcd(2n,k)\leq n$ cycles of $i\mapsto i+k$; if
$\sigma\rho(w)=\mathrm{rot}_k(w)$ then $w_{k-i}=1-w_i$, and the reflection
$i\mapsto k-i$ has no fixed points, so it pairs the $2n$ positions.
Either way at most $2^n$ words per $k$, fewer than $4n\cdot2^n$ words in
all, and each primitive class contains $2n$ words.  (b) is cited.

(vii) restates (K4), (K5), (i) and (iii).

(viii) With $h=\hat\lambda\equiv1\pmod4$, $h^2+6h+1\equiv8\pmod{16}$, and
$p$ divides the same number, so $8p$ divides it and both $b$ and $m$ are
even while $-6-h$ and $h$ are odd: $-M\equiv I\pmod2$, and
$\operatorname{tr}(-M)=6$.  A positive word with four or more cyclic
runs has trace $\geq\operatorname{tr}((RL)^2)=7$, so a trace-$6$ word is
$R^aL^b$ with $2+ab=6$: one of $R^4L$, $R^2L^2$, $RL^4$.  Of these only
$R^2L^2=\begin{pmatrix}5&2\\2&1\end{pmatrix}$ is $\equiv I\pmod2$, a
conjugacy-invariant condition.  Hence $-M\sim R^2L^2$.  (Directly,
$X=(M+5I)/2$ is integral, commutes with $M$ and has determinant
$(1-30+25)/4=-1$.)  Chirality of $L(p,J)$: $J^2\equiv-1$ with
$2J^2\equiv1$ would force $p\mid3$.
\end{proof}

\begin{remark}
The class $R^2L^2$ contains, for every admissible $(J,p)$ and both
$m\in\{2p,4p\}$, the matrix $-M$ of Lemma~\ref{lem:R}; this is the
precise sense in which the zero of Theorem~\ref{thm:C3} is a
``witnessed'' zero on the Sol side while, for $p\neq7$, having no
orientation-reversing or level-balance mechanism on the lens-space side
(at $p=7$ it is level-wise balanced, Proposition~\ref{prop:C3-balance}).
Neither statement
transfers symmetries between the two geometries; the zero itself
cannot identify why it is zero.
\end{remark}

\section{The silver-rank zero set}
\label{sec:silver-zero-set}

At a silver rank $\eta_J$ is an even integer at every modulus
(Corollary~\ref{cor:E}(b) and Theorem~\ref{thm:denominator}(c)), and its
zeros include the balanced family (b) of
Theorem~\ref{thm:classification}, the orbit $q\equiv\pm J,\pm2J$ of
Theorem~\ref{thm:C3}, and the mirror orbit when $-1$ is a square modulo
$p$.  At the level $p=2J^2-1$ itself the zero set is the subject of the
companion paper \cite{SilverRankZeros2026}.  We record its main statements without
proofs; the numbering refers to \cite{SilverRankZeros2026}, and nothing
else in this paper depends on them.

\begin{itemize}
\item \emph{The $S$-unit coordinate} \cite[Lem.~4.1,
Prop.~4.3]{SilverRankZeros2026}.  Reading $\Z[\sqrt2]$ modulo $p$
through $\sqrt2\mapsto2J$, the four zero families of
\cite[Thm.~7.1]{MahmoodP1Collection2026} are the orbits of $\sqrt2$,
$\delta/\sqrt2$, $\delta$ and $2\delta$, $\delta=1+\sqrt2$.  In the
Cayley coordinate $\rho=(1+q)/(1-q)$ the family (b), the orbit of
Theorem~\ref{thm:C3} and the mirror orbit are $\rho\in\{\pm\delta^{\pm1}\}$,
$\rho\in\{\pm\delta^{\pm2}\}$ and $\rho^2\equiv-1$.
\item \emph{A fifth family} \cite[Thm.~5.1]{SilverRankZeros2026}:
$\eta_J(L(2J^2-1,\,3-4J))=(J-1)\ind_{J\text{ odd}}$.
\item \emph{The Pell ladder} \cite[Thm.~6.1]{SilverRankZeros2026}.  With
$\delta^n=a_n+P_n\sqrt2$ and $q_n=a_n+2P_nJ$: if $P_n\mid J$ then
$\eta_J(L(p,q_n))=0$ for odd $n$ and $\frac23P_{n-2}J$ for even $n$; if
$n\geq2$ and $J\equiv1\pmod{P_n}$ then $\eta_J(L(p,q_n))=\gamma_n(J-1)/P_n$
for an explicit $\gamma_n>0$.
\item \emph{The period law} \cite[Conj.~7.1, Thm.~8.4,
Cor.~8.5]{SilverRankZeros2026}.  Conjecturally every nonconstant $S$-unit orbit is
eventually quasi-polynomial in $J$ with period dividing its coefficient
of $\sqrt2$, supported by a reported out-of-sample test on $64$ orbits
(only summary counts were retained; see the companion's Appendix~A).
This is proved for every nonconstant orbit of norm $\pm1$ or
$\pm2$ (the powers $\delta^n$, $n\geq1$, and the norm-two orbits
$a_nJ\pm P_n$), with explicit thresholds past which $\eta_J$ is linear in
$J$ on each residue class, complete zero sets for $24$ families, and
eventually periodic zero sets; the constant orbit $q=\pm1$ lies outside
the law, and there $\eta_J(L(2J^2-1,1))=-J(J-1)(2J-1)/3$ is cubic.
\item \emph{Rank-uniform zeros} \cite[Cor.~8.9,
Cor.~7.2(a)]{SilverRankZeros2026}.  Among the orbits of norm $\pm1$ and
$\pm2$ exactly $q\equiv J$, $J+1$, $2J-1$ vanish at every rank; in a
window of $64$ orbits exactly the four families of
\cite{MahmoodP1Collection2026} do.
\item \emph{Unbounded ambiguity} \cite[Cor.~8.13]{SilverRankZeros2026}.
For every $K$ there is a rank $J$ at which $K$ pairwise non-homeomorphic
chiral lens spaces $L(2J^2-1,q_i)$, none rank-$J$ balanced
(Theorem~\ref{thm:classification}), all have $\eta_J=0$.
\end{itemize}

The last item is the sharp form of the gap between vanishing and balance
of Example~\ref{ex:zero-not-balanced}.  By Corollary~\ref{cor:rank-one}
a zero of $\eta_1$ is exactly a mirror symmetry, and by
Corollary~\ref{cor:E}(c) the same holds at prime moduli off the silver
ranks with $p\nmid J$; at silver ranks $J\geq2$ a single vanishing $\eta_J$ can hide an
unbounded amount of handedness.  The companion also records a refuted
conjecture \cite[Conj.~4.6]{SilverRankZeros2026} and the reason the
Rademacher mechanism of Section~\ref{sec:silver} does not close the
classification \cite[Rem.~4.5]{SilverRankZeros2026}.

\section{Open problems}
\label{sec:open}

\begin{enumerate}[label=(\arabic*)]
\item \emph{Silver-rank zeros.}  Determine, for $p\mid2J^2-1$, the set of
$q$ with $\eta_J(L(p,q))=0$ (Section~\ref{sec:silver-zero-set} and
\cite[Sec.~9]{SilverRankZeros2026}); in particular prove the period law
\cite[Conj.~7.1]{SilverRankZeros2026} beyond norm $\pm2$, and describe
the zeros outside the $S$-unit sector.
\item \emph{Composite moduli.}  Theorem~\ref{thm:denominator}(b),(c) gives
necessary local conditions for a zero; they are not sufficient
($\eta_{25}(L(25,1))=-104$).  Observation~\ref{obs:composite} shows zeros
in which several nonzero primitive parts cancel.  A classification of
the zeros of $\eta_J$ at composite $p$, even off the silver ranks of the
prime factors, is open.
\item \emph{The prime $3$ and $p=2,3$ boundaries.}  Theorem~\ref{thm:E}
fails at $p=3$ ($\eta_1(L(3,1))=-2/9$) and Theorem~\ref{thm:denominator}(d)
needs $3\nmid p$; the $3$-primary behaviour of $\eta_J$ is described only
by Theorem~\ref{thm:denominator}(c).
\item \emph{Literature.}  Congruences of the type of Theorems~\ref{thm:E}
and \ref{thm:U} may be instances of index-theoretic congruences for
flat bundles; and Corollary~\ref{cor:S} is likely classical.  Neither
has been searched for.
\item \emph{Other operators.}  Whether the level-wise criterion of
Theorem~\ref{thm:criterion} has an analogue for twisted or half-integer
spin operators on lens spaces, where the eta invariants are known
\cite{Bar1996,Boldt2017}, is not addressed here.
\end{enumerate}

\appendix

\section{Exact computations}
\label{app:verification}

Every statement in this paper was accompanied by exact computations,
using integer or rational arithmetic (or certified interval arithmetic
where a transcendental identity is involved), and in several cases by
predictions fixed before the corresponding rows were computed
(``held-out'' below; ``opened'' rows were available when the statement
was formulated).  These computations test the identities and the
conventions; they are not substitutes for the proofs.  Numerical values
quoted as ``worst difference'' are high-precision floating-point
corroboration, not interval enclosures.

\begin{small}
\begin{longtable}{@{}L{0.24\textwidth}L{0.41\textwidth}L{0.2\textwidth}L{0.07\textwidth}@{}}
\toprule
statement & what was checked & rows & failures\\
\midrule
\endhead
Branch counts \eqref{eq:branch-counts} & direct weight-pair enumeration
against a Clebsch--Gordan recurrence on residue-class character vectors,
$p\leq40$, all $q$, $J\in\{1,2\}$, levels $\leq2p+10$ & $63{,}198$
(opened) & $0$\\
Eta anchor & the regularized level sum equals $\eta_J$ from the master
formula \eqref{eq:master} as emitted by an independent certified
implementation, $p\leq30$, all $q$, $J\in\{1,2,3\}$ & $831$ (opened) & $0$\\
Corollary~\ref{cor:affine} & periodic increments of residue-class weight
counts, $p\leq60$, every multiplier; windows $n<2p$ and $n<6p$ give
identical verdicts, $p\leq40$, $J\leq4$ & all rows & $0$\\
Corollary~\ref{cor:rank-one} & rank-one balance $\iff q^2\equiv-1$,
$p\leq60$; $s(h,k)=0\iff h^2\equiv-1$, $k\leq300$ & all rows & $0$\\
Rank $4$ prediction & chiral balanced spaces at $J=4$, $p\leq64$, predicted
as $p\in\{31,62\}$ with four $q$ each before computation & held-out & hit\\
Theorem~\ref{thm:kernel} & \eqref{eq:kernel-form} evaluated in $200$-bit
interval arithmetic contains the exact integer $pD_J(n)$, $p\leq30$, all
$q$, $J\leq8$, $n\leq2p$ & $91{,}840$ (opened) & $0$\\
Theorem~\ref{thm:criterion} & criterion (i)$\wedge$(ii) equals direct
balance on $n<6p$, $p\leq260$, all $q$, $J\leq8$; the symbolic identities
of the proof verified in a computer algebra system & $164{,}784$ & $0$\\
Example~\ref{ex:balanced-table} & chiral balanced sets for $J\leq8$,
$p\leq260$, and eight rank scans, predicted from the families before the
rows $J\geq5$ were computed & held-out & hit\\
Theorem~\ref{thm:classification} (support) & criterion $\iff$ (a)$\vee$(b)$\vee$(c)
for all $p\leq300$, all $q$, all $J$ mod $p$ & $5{,}483{,}516$ & $0$\\
Lemma~\ref{lem:parity-projection} & balance of $L(2p,\tilde q)$ equals
balance of $L(p,q)$ by the criterion, odd $p\leq151$, all $q$, $J\leq2p$
& $954{,}920$ & $0$\\
Remark~\ref{rem:certificate} & class-generator membership over $\Z$ in
the unit-saturated ideal, tested for each of the $105$ matchings (one
representative basis per class retained); semantic check on every odd
$p\in[5,101]$ and the
composites $121$, $125$, $143$, $169$, $175$, $187$, $221$, $245$, $289$,
$323$, $343$, $385$ (all units
$q$, all $J$ mod $p$); even reduction for $p\equiv2\pmod4$, $p\leq402$
& $105$; $712{,}060$; $1{,}116{,}462$ & $0$\\
Theorem~\ref{thm:decomposition} & primitive parts by the element route
\eqref{eq:generic-element}--\eqref{eq:clifford-element} against M\"obius
inversion, $d\leq60$, all $q$, $J\leq8$ (worst difference
$1.7\cdot10^{-55}$ at $60$ digits) & $8{,}808$ & $0$\\
Proposition~\ref{prop:first-moment} & $\sum_ru_r=0$ and
$\sum_r[v_r-a_ru_r]=0$, $p<40$, all $q$, $J\in\{1,2,3,5\}$ & $1{,}888$ & $0$\\
Observation~\ref{obs:prime-zeros} & zeros of $\eta_J$ on primes
$41$--$199$ ($J\leq12$) and $41$--$101$ ($13\leq J\leq(p-1)/2$); the
prediction, fixed in advance, ``zeros are mirror, silver-family, or the four families
at $p=2J'^2-1$'' was refuted by $196$ zeros in $49$ orbits & $71{,}230$
(held-out) & $196$ unpredicted\\
Observation~\ref{obs:composite} & zeros on composites $41$--$120$,
$J\leq10$, classified by primitive parts & $25{,}180$ & --\\
Observation~\ref{obs:silver-blind} & dichotomy, silver integrality and
$q=J$ on primes $211$--$599$ never inspected before the statements were
fixed & $67{,}032+41{,}286+12{,}210$ (held-out) & $0$\\
Theorem~\ref{thm:E} & symbolic Laurent coefficients of $f$, $f_{J,1}$ and
the trace constants (two independent computer algebra systems); trace
form \eqref{eq:trace-form} against $p\,\eta_J$, primes $5$--$47$, all $q$,
$J\leq6$ (worst difference $8.7\cdot10^{-57}$); (i), (ii),
Corollary~\ref{cor:E}(a) on all $17$ primes $601$--$701$, all $q$,
$J\in\{1,2,3\}$, plus silver ranks $J_s,J_s+1$ and $J=p$ at $601,607$;
master formula against the evaluator, $p\leq40$, $J\leq8$; Rademacher's
congruence, primes $5$--$199$ & $1{,}860$; $43{,}026$ (held-out); $3{,}912$;
$4{,}178$ & $0$\\
Corollary~\ref{cor:E}(d) & $\eta_{np}(L(p,\pm1))=\mp n(p^2-1)/6$ at
$p\in\{5,7,11,13\}$, $n\in\{1,2\}$; $\eta_J-J(2J^2-1)\eta_1\in2\Z$ on
primes $5$--$59$, all $q$, $J\leq12$ & $8$; $5{,}040$ & $0$\\
Theorem~\ref{thm:U} & master formula against the evaluator,
$41\leq p\leq60$, $J\leq8$; Lemma~\ref{lem:dedekind-facts}, $p\leq400$;
Lemma~\ref{lem:low-mode-parity}, $p\leq80$, $J\leq16$; \eqref{eq:U} on
every composite $121$--$220$ ($J\leq6$), on $p\in\{2,3\}$
($9\leq J\leq60$) and on $p\in\{4,8,9,16,25,27,32,49\}$ ($9\leq J\leq24$);
Theorem~\ref{thm:denominator}(b) on the $364$ zeros found & $4{,}896$;
$48{,}677$; $31{,}440$; $45{,}072+156+1{,}856$ (held-out except $p=121,150,200$
and $p\in\{2,3\}$, which were computed earlier) & $0$\\
Theorem~\ref{thm:C3} & the chain $K=0$, $16K=T-p$, parity split,
Lemma~\ref{lem:D} at $m\in\{2p,4p\}$, Lemma~\ref{lem:R} by two
Dedekind-sum routes, $\Psi(M)=0$ and content $2$, on every divisor
$p>1500$ of $2J^2-1$ for $60\leq J\leq110$ and at both silver ranks of the
first $12$ primes $\equiv1$ and $12$ primes $\equiv7\pmod8$ above
$100{,}000$ (held-out), and on $p\leq1500$, $60\leq J\leq110$ (opened);
$\eta_J(L(p,J))=2K/p^2=0$ by the exact evaluator for $p\leq12000$;
Corollary~\ref{cor:S} on every row; family (b) on $15$ rows &
$62+48$ (held-out) $+73$; $29$; $183$; $15$ & $0$\\
Theorem~\ref{thm:Z} & (K4), \eqref{eq:psi-reflection}, (K2), word route
against the Gauss reduced-form route, inert $\iff$ negative Pell, on every
primitive word of length $\leq18$; $M_1$, $M_2$; all primitive balanced
classes of length $\leq8$; two-block $a,b\leq60$; primitive four-block
tuples with $a+b+c+d\leq32$ ($1{,}232$ zeros, $1{,}008$ accidental);
the family $(2,1,n-1,n)$, $n\leq60$; $B_{2n}$ and $W_{2n}$ at every even
length $\leq24$ ($B_{22}=32{,}065$, $W_{22}=1{,}085$; $B_{24}=112{,}632$,
$W_{24}=2{,}150$); $-M\sim R^2L^2$ for $J\leq40$, all divisors, both $m$ &
$31{,}040$; $13$; $3{,}600$; $35{,}840$; $58$; $12$ lengths; $154$ & $0$\\
Section~\ref{sec:silver-zero-set} & the exact computations behind the
companion's statements (fifth family, Pell ladder, period-law test and
proof, complete zero sets, the refuted Conjecture~4.6) are listed in
\cite[App.~A]{SilverRankZeros2026} & -- & --\\
\bottomrule
\end{longtable}
\end{small}

Share of primitive balanced classes with an inert or reciprocal witness,
$W_{2n}/B_{2n}$: $6/8$ at length $8$, $29/75$ at $12$, $128/800$ at $16$,
$531/9{,}225$ at $20$, $1{,}085/32{,}065$ at $22$, $2{,}150/112{,}632$ at $24$.

In each block, deliberately mutated versions of the statements (a wrong
sign, a dropped term, a wrong root, a wrong residue class) were run and
failed at a named witness; these mutation checks guard against a checker
that passes vacuously.  One refutation of a prediction fixed in advance
is recorded above: the conjectured prime-modulus zero list was false,
and the zeros it missed led to Observation~\ref{obs:prime-zeros} and
thence to Theorems~\ref{thm:E} and \ref{thm:C3}.  Two further
refutations, of the companion's Conjecture~4.6 at $p=1129$ and of a control premise at
$q=4-4J$, are recorded in \cite[App.~A]{SilverRankZeros2026}.

\section*{Authorship, assistance and rights}
Luqman Mahmood directed this work and is responsible for its claims and final
text. OpenAI Codex and Anthropic Claude assisted with mathematical development,
computation, checking and editing. Internal and automated checks are not
external peer review. Computational tests do not replace proofs of universal
statements, and no physical identification or experiment is implied.
The paper and mathematical data are licensed CC BY 4.0. Software source and
executable mathematical verification scripts are excluded from that grant;
their rights are reserved unless explicitly stated otherwise.

\begin{thebibliography}{99}

\bibitem{MahmoodP1Collection2026}
L. Mahmood, \emph{Eta invariants of transverse-traceless tensor curl on
three-dimensional lens spaces}, companion preprint, 2026,
\texttt{doi:10.5281/zenodo.21433675}.

\bibitem{SilverRankZeros2026}
L. Mahmood, \emph{Silver-rank zeros of the tensor-curl eta invariant on
lens spaces: silver $S$-units, a Pell ladder, and a period law},
companion preprint, 2026. DOI:
\url{https://doi.org/10.5281/zenodo.23049615}.

\bibitem{MahmoodP3Collection2026}
L. Mahmood, \emph{Dirac-to-tensor-curl eta transfer on round spherical
space forms with affine-ADE specializations}, companion preprint, 2026,
\texttt{doi:10.5281/zenodo.21433687}.

\bibitem{MahmoodP7Collection2026}
L. Mahmood, \emph{Square elliptic completions and APS boundary
information on round spherical quotients}, companion preprint, 2026,
\texttt{doi:10.5281/zenodo.22917346}.

\bibitem{APS1975}
M. F. Atiyah, V. K. Patodi and I. M. Singer, Spectral asymmetry and
Riemannian geometry. II, \emph{Math. Proc. Cambridge Philos. Soc.}
\textbf{78} (1975), 405--432.

\bibitem{Atiyah1987}
M. F. Atiyah, The logarithm of the Dedekind $\eta$-function,
\emph{Math. Ann.} \textbf{278} (1987), 335--380.

\bibitem{Bar1996}
C. B\"ar, The Dirac operator on space forms of positive curvature,
\emph{J. Math. Soc. Japan} \textbf{48} (1996), 69--83.

\bibitem{Bar2019}
C. B\"ar, The curl operator on odd-dimensional manifolds,
\emph{J. Math. Phys.} \textbf{60} (2019), 031501; arXiv:1702.02044.

\bibitem{Basu2015}
R. Basu, Higher spin de Sitter quantum gravity, \emph{J. High Energy
Phys.} \textbf{2015}, no.~10, 151; arXiv:1507.04757.

\bibitem{BeckRobins2004}
M. Beck and S. Robins, Dedekind sums: a combinatorial-geometric
viewpoint, in \emph{Unusual applications of number theory}, DIMACS Ser.
Discrete Math. Theoret. Comput. Sci. \textbf{64}, Amer. Math. Soc., 2004,
25--35; arXiv:math/0112076.

\bibitem{Boldt2017}
S. Boldt, Properties of the Dirac spectrum on three dimensional lens
spaces, \emph{Osaka J. Math.} \textbf{54} (2017), 747--765; arXiv:1504.03121.

\bibitem{Burrin2022}
C. Burrin, The Manin--Drinfeld theorem and the rationality of Rademacher
symbols, \emph{J. Th\'eor. Nombres Bordeaux} \textbf{34} (2022), 739--753.

\bibitem{BurrinVonEssen2024}
C. Burrin and F. von Essen, Windings of prime geodesics, \emph{Int.
Math. Res. Not. IMRN} \textbf{2024}, no.~22, 13931--13963,
\texttt{doi:10.1093/imrn/rnae225}; arXiv:2209.06233.

\bibitem{CastroLashkariMaloney2011Farey}
A. Castro, N. Lashkari and A. Maloney, A de Sitter Farey tail,
\emph{Phys. Rev. D} \textbf{83} (2011), 124027; arXiv:1103.4620.

\bibitem{COMPACTIIIa2024}
COMPACT Collaboration (A. Samandar et al.), Cosmic topology. Part IIIa.
Microwave background parity violation without parity-violating
microphysics, \emph{J. Cosmol. Astropart. Phys.} \textbf{11} (2024), 020;
arXiv:2407.09400.

\bibitem{Ghys2006}
\'E. Ghys, Knots and dynamics, in \emph{Proceedings of the International
Congress of Mathematicians (Madrid, 2006)}, Vol.~I, European
Mathematical Society, Z\"urich, 2007, 247--277.

\bibitem{Girstmair2016}
K. Girstmair, On the values of Dedekind sums, arXiv:1610.05588.

\bibitem{Jeffrey1992}
L. C. Jeffrey, Chern--Simons--Witten invariants of lens spaces and torus
bundles, and the semiclassical approximation, \emph{Comm. Math. Phys.}
\textbf{147} (1992), 563--604.

\bibitem{Katase1987}
K. Katase, On the value of Dedekind sums and eta-invariants for
3-dimensional lens spaces, \emph{Tokyo J. Math.} \textbf{10} (1987),
327--347.

\bibitem{MartinSvesko2020}
V. L. Martin and A. Svesko, Higher spin partition functions via the
quasinormal mode method in de Sitter quantum gravity, \emph{SciPost Phys.}
\textbf{9} (2020), 039; arXiv:2004.00128.

\bibitem{Matsusaka2024}
T. Matsusaka, A hyperbolic analogue of the Rademacher symbol,
\emph{Math. Ann.} \textbf{388} (2024), 2843--2886;
\href{https://doi.org/10.1007/s00208-023-02588-9}{doi:10.1007/s00208-023-02588-9};
arXiv:2003.12354v4.

\bibitem{Meyer2005}
J. L. Meyer, Symmetric arguments in the Dedekind sum, \emph{Fibonacci
Quart.} \textbf{43} (2005), 122--123.

\bibitem{PrzytyckiYasuhara2003}
J. H. Przytycki and A. Yasuhara, Symmetry of links and classification of
lens spaces, \emph{Geom. Dedicata} \textbf{98} (2003), 57--61;
arXiv:math/0011119.

\bibitem{RademacherGrosswald1972}
H. Rademacher and E. Grosswald, \emph{Dedekind sums}, Carus Mathematical
Monographs \textbf{16}, Mathematical Association of America, 1972.

\bibitem{Sarnak2007}
P. Sarnak, Reciprocal geodesics, \emph{Clay Math. Proc.} \textbf{7}
(2007), 217--237.

\bibitem{SunWangWu2008}
H. Sun, S. Wang and J. Wu, Self-mapping degrees of torus bundles and
torus semi-bundles, \emph{Osaka J. Math.} \textbf{47} (2010), 131--155;
\href{https://doi.org/10.18910/4707}{doi:10.18910/4707}; arXiv:0810.1799.

\bibitem{Zagier1975}
D. Zagier, Nombres de classes et fractions continues, in
\emph{Journ\'ees Arithm\'etiques de Bordeaux (1974)},
\emph{Ast\'erisque} \textbf{24--25} (1975), 81--97.

\end{thebibliography}

\end{document}
